\documentclass[UTF8, 12pt, fontset=fandol, oneside]{ctexart}
\usepackage{researchnotes}

% 保留分册独立编译时的章号前缀。
\renewcommand{\thesection}{5.\arabic{section}}

\title{第五章\quad 导数的应用}
\author{}
\date{}

\begin{document}

\maketitle

上一章中, 我们着重介绍了导数的概念及计算导数的各种方法. 这一章, 我们将主要利用导数来研究函数的性质。

\section{微分中值定理}

微分中值定理是反映函数和导数之间联系的重要定理, 是利用导数研究函数的桥梁, 因此它是一个强有力的工具. 灵活地运用它可使许许多多问题的解决变得轻而易举. 在建立微分中值定理时需要用到如下的费马 (Fermat) 定理, 因此我们先来介绍这一定理。

\subsection{费马定理}

极值问题是数学中的重要研究对象, 在本章中我们将利用导数来讨论此类问题. 为此我们引入

\noindent\textbf{定义 5.1.1}\quad 若函数 $f(x)$ 在 $x_0$ 的某个去心邻域 $U_0(x_0,\delta)$ 内恒有
\[
  f(x)\le f(x_0)\quad (f(x)\ge f(x_0)),
\]
则称 $x_0$ 为 $f(x)$ 的极大 (小) 值点, $f(x_0)$ 称为它的极大 (小) 值; 若在上述定义中, 不等号严格成立, 则称 $f(x_0)$ 为严格极大 (小) 值. 极大值点和极小值点统称为极值点, 极大值和极小值统称为极值。

从图像上来看, 一个函数 $y=f(x)$ 若在 $x_0$ 处取极值, 并且该函数的图像在点 $(x_0,f(x_0))$ 处存在切线, 则该切线一定是水平的. 事实上, 我们有

\noindent\textbf{定理 5.1.1(费马定理)}\quad 设函数 $f(x)$ 在 $U(x_0,\delta)$ 中有定义. 若 $x_0$ 是 $f(x)$ 的极值点, 且 $f'(x_0)$ 存在, 则 $f'(x_0)=0$。

\noindent\textbf{证明}\quad 无妨设 $f(x_0)$ 为极大值, 则当 $\Delta x>0$ 且 $x_0+\Delta x\in U(x_0,\delta)$ 时, 有
\[
  \frac{f(x_0+\Delta x)-f(x_0)}{\Delta x}\le0.
\]
令 $\Delta x\to0+0$, 得 $f'(x_0)\le0$。

当 $\Delta x<0$ 且 $x_0+\Delta x\in U(x_0,\delta)$ 时, 有
\[
  \frac{f(x_0+\Delta x)-f(x_0)}{\Delta x}\ge0.
\]
令 $\Delta x\to0-0$, 得 $f'(x_0)\ge0$。

综上所证, 推得 $f'(x_0)=0$。

\noindent\textbf{注}\quad 使得 $f'(x)=0$ 的点 $x_0$ 也称为 $f(x)$ 的驻点. 费马定理告诉我们, 若极值点可导, 则它必是驻点. 例子 $y=x^3$ 在 $x=0$ 处的性质告诉我们, 驻点未必是极值点。

\subsection{罗尔微分中值定理}

\noindent\textbf{定理 5.1.2(罗尔 (Rolle) 微分中值定理)}\quad 设函数 $f(x)$ 在 $[a,b]$ 上连续, 在 $(a,b)$ 内可导, 且 $f(a)=f(b)$, 则在 $(a,b)$ 内至少存在一点 $\xi$, 使得
\[
  f'(\xi)=0.
\]

\noindent\textbf{证明}\quad 因为 $f(x)$ 在 $[a,b]$ 连续, 所以 $f(x)$ 在 $[a,b]$ 上有最大值 $M$ 与最小值 $m$. 如果 $M=m$, 则 $f(x)$ 在 $[a,b]$ 为一常值函数, 从而 $f'(x)=0,\ \forall x\in[a,b]$. 此时, 我们可取 $(a,b)$ 中任意一点作为 $\xi$. 如果 $M>m$, 则 $M$ 与 $m$ 至少有一个不等于 $f(a)=f(b)$, 不妨设 $M>f(a)$. 这时, 必存在 $\xi\in(a,b)$, 使得 $f(\xi)=M$, 即 $\xi$ 为 $f(x)$ 的一个极值点. 再由 $f'(\xi)$ 存在, 并且应用费马定理, 得 $f'(\xi)=0$。

罗尔微分中值定理的几何意义是, 在一段每点都有切线的曲线上, 若两端点的高度相同, 则在此曲线上至少存在一条水平切线 (参见图 5.1.1)。

此外, 罗尔微分中值定理只是告诉我们在区间 $(a,b)$ 内存在一点 $\xi$ 使得 $f'(\xi)=0$, 但并没有告诉我们怎样去求这一 $\xi$. 事实上, 一般来讲, 要将这一 $\xi$ 求出来是十分困难. 其次, 在很多时候我们会发现, 满足这样要求的 $\xi$ 可能有很多, 甚至有可能是无穷多个。

\noindent\textbf{例 5.1.1}\quad 设函数 $f(x)$ 在 $(a,b)$ 上可微, 证明 $f(x)$ 在 $(a,b)$ 内的两个零点之间必有 $f(x)+f'(x)$ 的零点。

\noindent\textbf{证明}\quad 令 $F(x)=e^x f(x)$, 由罗尔微分中值定理知, 在 $f(x)$ 的两个零点之间必有
\[
  F'(x)=e^x(f(x)+f'(x))
\]
的零点, 也即有 $f(x)+f'(x)$ 的零点。

\noindent\textbf{例 5.1.2}\quad 设 $P_n(x)$ 为一个 $n$ 次多项式, 证明 $e^x-P_n(x)=0$ 至多有 $n+1$ 个不同的根。

\noindent\textbf{证明}\quad 令 $f(x)=e^x-P_n(x)$. 假若 $f(x)$ 有 $n+2$ 个不同的零点, 则由罗尔微分中值定理知, $f'(x)=e^x-(P_n(x))'$ 至少有 $n+1$ 个不同的零点; 然后再应用罗尔微分中值定理, 又有 $f''(x)=e^x-(P_n(x))''$ 至少有 $n$ 个不同的零点. 如此下去, 即知 $f^{(n+1)}=e^x-(P_n(x))^{(n+1)}$ 至少有一个零点. 由于 $P_n(x)$ 是一个 $n$ 次多项式, 所以这就蕴涵着 $e^x$ 必有一个零点, 这显然是不可能的. 这一矛盾说明 $e^x-P_n(x)=0$ 至多有 $n+1$ 个不同的根。

\noindent\textbf{例 5.1.3}\quad 设函数 $f(x)$ 在 $[0,1]$ 上连续, 在 $(0,1)$ 内可导, 且 $f(0)=f(1)=0$. 再假设 $f(t_0)=\alpha$, 其中 $t_0\in(0,1)$. 证明: 存在 $\xi\in(0,1)$, 使得 $f'(\xi)=\alpha$。

\noindent\textbf{证明}\quad 若 $\alpha=0$, 利用罗尔微分中值定理即得。

现设 $\alpha\ne0$. 构造辅助函数
\[
  F(x)=f(x)-\alpha x,
\]
则 $F(x)$ 在 $[0,1]$ 上连续, 在 $(0,1)$ 内可导, 并且有
\[
  F(0)=0,\quad F(1)=-\alpha,\quad F(t_0)=(1-t_0)\alpha.
\]
由于 $F(1)F(t_0)<0$, 因此在 $(t_0,1)$ 中必存在 $\eta$, 使得 $F(\eta)=0$. 因此我们在 $[0,\eta]$ 上对 $F(x)$ 应用罗尔微分中值定理, 可知在 $(0,\eta)\subset(0,1)$ 中存在 $\xi$, 使得 $F'(\xi)=0$, 即 $f'(\xi)=\alpha$。

\subsection{拉格朗日微分中值定理}

罗尔微分中值定理要求所考虑的函数在所讨论区间的两端点有相同的函数值, 这给它的应用范围带来很大的限制. 因此, 一个自然的问题就是: 如果将罗尔微分中值定理中的条件 ``$f(a)=f(b)$'' 去掉, 将有什么结论呢? 其实, 如果将坐标轴去掉, 纯粹从几何上来看, 罗尔微分中值定理是说, 在一段每点都有切线的曲线上至少有一点的切线平行于两端点的连线. 而这正是一个可微函数的图像所具有的几何特征. 因此, 如果将罗尔微分中值定理中的条件 ``$f(a)=f(b)$'' 去掉, 则有如下的定理:

\noindent\textbf{定理 5.1.3(拉格朗日 (Lagrange) 微分中值定理)}\quad 若函数 $f(x)$ 在 $[a,b]$ 上连续, 在 $(a,b)$ 内可导, 则在 $(a,b)$ 内至少存在一点 $\xi$, 使得
\begin{equation}
  f'(\xi)=\frac{f(b)-f(a)}{b-a}.
  \tag{5.1.1}
\end{equation}

\noindent\textbf{证明}\quad 考查 $f(x)$ 减去过 $(a,f(a))$ 和 $(b,f(b))$ 这两点的直线所确定的线性函数而得到的函数:
\[
  F(x)=f(x)-\left[f(a)+\frac{f(b)-f(a)}{b-a}(x-a)\right].
\]
容易验证, 这样定义的函数 $F(x)$ 满足罗尔微分中值定理的所有条件. 因此, 存在 $\xi\in(a,b)$, 使得
\[
  F'(\xi)=f'(\xi)-\frac{f(b)-f(a)}{b-a}=0,
\]
即
\[
  f'(\xi)=\frac{f(b)-f(a)}{b-a}.
\]

\noindent\textbf{注}\quad (1) 在上述定理中, 若我们只假定 $f(x)$ 在以 $a$ 和 $b$ 为端点的闭区间满足所有的条件, 不难看出, 对 $a>b$ 的情形, 拉格朗日微分中值定理仍具有同样的形式。

(2) 公式 (5.1.1) 有时也称为拉格朗日中值公式. 在实际应用时, 常将拉格朗日中值公式写成如下形式
\[
  f(a+h)-f(a)=f'(\xi)h,
\]
这里 $h=b-a$ 可以是正数, 也可以是负数, 而 $\xi$ 则位于 $a$ 与 $a+h$ 之间. 有时, 为了避免 ``$\xi$ 位于 $a$ 与 $a+h$ 之间'' 可能出现的两种情形, 常将上式改写成
\[
  f(a+h)=f(a)+f'(a+\theta h)h,
\]
其中 $\theta=\dfrac{\xi-a}{h}$. 无论 $h$ 是正数还是负数, 总有 $0<\theta<1$, 这是因为 $\xi$ 位于 $a$ 与 $a+h$ 之间的缘故。

此外, 拉格朗日微分中值定理的几何意义如图 5.1.2 所示, 即在一段每点都有切线的曲线上至少有一点的切线平行于两个端点的连线。

由拉格朗日微分中值定理立即可推出如下两个十分有用的结论:

\noindent\textbf{推论 1}\quad 设函数 $f(x)\in C[a,b]$, 且在 $(a,b)$ 内可导. 若 $f'(x)=0$, 则 $f(x)\equiv C$, 其中 $C$ 是一常数。

\noindent\textbf{证明}\quad 对 $\forall x\in[a,b]$, 在 $[a,x]$ 上应用拉格朗日微分中值定理知, 存在 $\xi\in(a,x)$, 使得
\[
  0=f'(\xi)=\frac{f(x)-f(a)}{x-a},
\]
由此即得
\[
  f(x)\equiv f(a)=C.
\]
推论 1 的结论似乎是显然的, 但如果不利用微分中值定理, 似乎很难证明它。

\noindent\textbf{推论 2}\quad 设函数 $f(x),g(x)\in C[a,b]$, 且在 $(a,b)$ 内可导. 若 $f'(x)=g'(x)$, 则 $f(x)=g(x)+C$, 其中 $C$ 是一常数。

\noindent\textbf{证明}\quad 对 $f(x)-g(x)$ 应用推论 1 即得。

\noindent\textbf{例 5.1.4}\quad 证明当 $x>-1$ 且 $x\ne0$ 时, 有
\[
  \frac{x}{1+x}<\ln(1+x)<x.
\]

\noindent\textbf{证明}\quad 对 $\ln(1+x)$, 在 $0$ 和 $x$ 之间应用拉格朗日微分中值定理, 得
\[
  \frac{\ln(1+x)}{x}
  =
  \frac{\ln(1+x)-\ln1}{x-0}
  =
  \left.(\ln(1+x))'\right|_{x=\xi}
  =
  \frac1{1+\xi},
\]
其中 $\xi$ 在 $0$ 和 $x$ 之间. 再注意到, 当 $x>0$ 时, 有
\[
  \frac1{1+x}<\frac1{1+\xi}<1;
\]
而当 $-1<x<0$ 时, 有
\[
  1<\frac1{1+\xi}<\frac1{1+x},
\]
便知所证不等式成立。

\noindent\textbf{例 5.1.5}\quad 证明 $f(x)=x^\alpha\ (0<\alpha<1)$ 在 $[0,+\infty)$ 上一致连续。

\noindent\textbf{证明}\quad 由于 $f(x)$ 在 $[0,+\infty)$ 上连续, 因此它在 $[0,2]$ 上一致连续. 于是, $\forall\varepsilon>0,\exists\delta_1>0$, 使得当 $x_1,x_2\in[0,2]$ 且 $|x_1-x_2|<\delta_1$ 时, 有
\[
  |f(x_1)-f(x_2)|<\varepsilon.
\]
另一方面, 因为 $f'(x)=\alpha x^{\alpha-1}$ 在 $[1,+\infty)$ 上严格单调递减, 所以在 $[1,+\infty)$ 上恒有 $|f'(x)|\le\alpha<1$. 于是, 对 $\forall x_1,x_2\in[1,+\infty)$, 应用拉格朗日微分中值定理, 得
\[
  |f(x_1)-f(x_2)|=|f'(\xi)(x_1-x_2)|\le |x_1-x_2|.
\]
这样, 对给定的 $\varepsilon$, 取 $\delta_2=\varepsilon$, 则当 $x_1,x_2\in[1,+\infty)$ 且 $|x_1-x_2|<\delta_2$ 时, 有
\[
  |f(x_1)-f(x_2)|<\varepsilon.
\]
现取 $\delta=\min\{\delta_1,\delta_2,1\}$, 则 $\forall x_1,x_2\in[0,+\infty)$, 当 $|x_1-x_2|<\delta$ 时, 一定有 $x_1,x_2\in[0,2]$ 或 $x_1,x_2\in[1,+\infty)$, 从而必有 $|f(x_1)-f(x_2)|<\varepsilon$. 这表明 $f(x)$ 在 $[0,+\infty)$ 上一致连续。

\noindent\textbf{注}\quad 从上例的证明中可以看出: 若一个函数在某区间上的导数有界, 则它必在该区间上一致连续. 此外, 上例也说明了, 一致连续的函数即使导数存在, 其导数也未必在所论区间上有界。

\noindent\textbf{思考题}\quad 当 $\alpha>1$ 时, 函数 $f(x)=x^\alpha$ 是否在 $[0,+\infty)$ 上一致连续?

\noindent\textbf{例 5.1.6}\quad 证明: 对任意的 $x\in\mathbb{R}$, 有
\[
  \arctan x=\arcsin\left(\frac{x}{\sqrt{1+x^2}}\right).
\]

\noindent\textbf{证明}\quad 令
\[
  f(x)=\arctan x-\arcsin\left(\frac{x}{\sqrt{1+x^2}}\right),
\]
则有
\[
  f'(x)
  =
  \frac1{1+x^2}
  -
  \frac1{\sqrt{1-\frac{x^2}{1+x^2}}}
  \cdot
  \frac{\sqrt{1+x^2}-\frac{x^2}{\sqrt{1+x^2}}}{1+x^2}
  =
  0,\quad \forall x\in\mathbb{R}.
\]
于是由推论 1 知 $f(x)\equiv C$, 而 $f(0)=0$, 故所证等式成立。

\subsection{柯西微分中值定理}

为了将拉格朗日微分中值定理作进一步的推广, 我们现在换个观点来看这个定理. 在该定理的条件下, 我们有 $f'(\xi_1)=\dfrac{f(b)-f(a)}{b-a}$, 其分子是函数 $f(x)$ 在区间端点之值的差, 而其分母可以看成函数 $y=x$ 在端点之值的差. 从这个观点来看, 一个自然的问题是: 能否将 $y=x$ 换成一个任意的函数呢? 下面的柯西微分中值定理回答了这个问题。

\noindent\textbf{定理 5.1.4(柯西微分中值定理)}\quad 若函数 $f(x)$ 和 $g(x)$ 在 $[a,b]$ 上连续, 在 $(a,b)$ 内可导, 而且 $g'(x)\ne0$, 则在 $(a,b)$ 内至少存在一点 $\xi$, 使得
\begin{equation}
  \frac{f(b)-f(a)}{g(b)-g(a)}=\frac{f'(\xi)}{g'(\xi)}.
  \tag{5.1.2}
\end{equation}

\noindent\textbf{注}\quad 分别应用拉格朗日微分中值定理于 $f(x)$ 和 $g(x)$, 我们可得
\[
  \frac{f(b)-f(a)}{g(b)-g(a)}=\frac{f'(\xi_1)}{g'(\xi_2)}.
\]
注意, 这里的 $\xi_1$ 与 $\xi_2$ 可能是不一样的! 因此, 我们不能简单地应用拉格朗日微分中值定理来证明本定理。

\noindent\textbf{证明}\quad 对函数 $g(x)$ 应用拉格朗日微分中值定理 (或罗尔微分中值定理), 可知 $g'(x)\ne0$ 蕴涵着 $g(a)\ne g(b)$. 仿照拉格朗日中值定理的证明方法, 令
\[
  G(x)=f(x)-\left[f(a)+\frac{f(b)-f(a)}{g(b)-g(a)}(g(x)-g(a))\right],
\]
则 $G(x)$ 满足罗尔定理的所有条件. 于是, 由罗尔定理知, 存在 $\xi\in(a,b)$, 使得 $G'(\xi)=0$. 整理即得公式 (5.1.2)。

柯西微分中值定理的几何意义是, 若 $uv$ 坐标平面的曲线由参数方程
\[
  \begin{cases}
    u=g(x),\\
    v=f(x),
  \end{cases}
  \quad x\in[a,b]
\]
给出, 其中 $f(x)$ 和 $g(x)$ 满足定理的条件, 则存在 $\xi\in(a,b)$, 使得过点 $(g(\xi),f(\xi))$ 的切线平行于两端点的连线。

此外, 类似于拉格朗日微分中值定理, 柯西微分中值定理也常写成如下形式:
\[
  \frac{f(b)-f(a)}{g(b)-g(a)}
  =
  \frac{f'(a+\theta(b-a))}{g'(a+\theta(b-a))},
\]
其中 $0<\theta<1$, 且不要求 $b>a$。

\noindent\textbf{例 5.1.7}\quad 设函数 $f(x)$ 在区间 $(0,1]$ 上连续, 在区间 $(0,1)$ 内可导, 而且
\[
  \lim_{x\to0+0}\sqrt{x}f'(x)=\ell.
\]
证明 $f(x)$ 在区间 $(0,1]$ 上一致连续。

\noindent\textbf{证明}\quad 由 $\lim\limits_{x\to0+0}\sqrt{x}f'(x)=\ell$ 可知, 存在 $M>0$ 和 $0<c<1$, 使得
\[
  |\sqrt{x}f'(x)|\le M,\quad \forall x\in(0,c].
\]
现在取 $x_1,x_2\in(0,c]$, $x_1\ne x_2$, 应用柯西微分中值定理, 可得
\[
  \frac{f(x_1)-f(x_2)}{\sqrt{x_1}-\sqrt{x_2}}
  =
  \left.\frac{f'(x)}{(\sqrt{x})'}\right|_{x=\xi}
  =
  2\sqrt{\xi}f'(\xi),
\]
其中 $\xi$ 位于 $x_1$ 和 $x_2$ 之间. 于是, 有
\[
  |f(x_1)-f(x_2)|\le 2M|\sqrt{x_1}-\sqrt{x_2}|.
\]
当然, 上式对 $x_1=x_2$ 亦成立. 再注意到
\[
  |\sqrt{x_1}-\sqrt{x_2}|^2
  \le |\sqrt{x_1}-\sqrt{x_2}|\,|\sqrt{x_1}+\sqrt{x_2}|
  =
  |x_1-x_2|,
\]
对 $\forall\varepsilon>0$, 只需取 $\delta=\dfrac{\varepsilon^2}{(2M)^2}$, 则当 $x_1,x_2\in(0,c]$ 且 $|x_1-x_2|<\delta$ 时, 便有
\[
  |f(x_1)-f(x_2)|
  \le 2M|\sqrt{x_1}-\sqrt{x_2}|
  \le 2M|x_1-x_2|^{\frac12}
  <\varepsilon.
\]
这表明 $f(x)$ 在 $(0,c]$ 上一致连续. 又 $f(x)$ 在区间 $(0,1]$ 上连续蕴涵着 $f(x)$ 在区间 $[c,1]$ 上一致连续, 从而便有 $f(x)$ 在区间 $(0,1]$ 上一致连续。

从本节内容可以看出, 罗尔微分中值定理是一个基本的结果, 拉格朗日微分中值定理是罗尔微分中值定理的直接推广, 而柯西微分中值定理则是拉格朗日微分中值定理的一个重要推广. 它们各自适合不同的场合, 是我们利用导数研究函数的强有力工具。

\section{洛必达法则}

在建立求导法则和求导公式的过程中, 极限的理论和一些具体的极限起着决定性的作用. 而有了导数理论和求导公式之后, 又可利用它解决极限理论中某些不定式的极限问题。

大家已经知道, 如果我们已知
\[
  \lim_{x\to a}f(x)=A,\quad \lim_{x\to a}g(x)=B,
\]
这里 $A,B,a$ 可以是有限数或无穷大, 那么除去某些例外情形, 我们可以利用下面的极限运算法则去求它们的和、差、积、商及幂--指数函数式的极限:
\begin{enumerate}
  \item $\lim\limits_{x\to a} f(x)\pm g(x)=A\pm B$;
  \item $\lim\limits_{x\to a} f(x)\cdot g(x)=AB$;
  \item $\lim\limits_{x\to a}\dfrac{f(x)}{g(x)}=\dfrac AB\ (B\ne0)$;
  \item $\lim\limits_{x\to a} f(x)^{g(x)}=A^B$。
\end{enumerate}
对于以上各式, 所说的例外情况分别为:
\begin{enumerate}
  \item[(a)] 在 (1) 中出现两个同号无穷大相减;
  \item[(b)] 在 (2) 中 $A$ 与 $B$ 之一为 $0$, 另一为无穷大;
  \item[(c)] 在 (3) 中 $A$ 与 $B$ 同时为 $0$, 或者同为无穷大;
  \item[(d)] 在 (4) 中 $A=1,B=\infty$, 或者 $A=B=0$, 或者 $A=\infty,B=0$。
\end{enumerate}
这些例外情形所涉及的极限类型统称为\textbf{不定式}, 分别记为:
\[
  \infty-\infty,\quad 0\cdot\infty,\quad \frac00,\quad \frac\infty\infty,\quad 1^\infty,\quad 0^0,\quad \infty^0.
\]
所有这些不定式本质上只有一种: $\dfrac00$. 其他的不定式经过简单的变换可以化为这种不定式, 如 $0\cdot\infty=0\cdot\dfrac10$ 就成为 $\dfrac00$ 型, 或者 $0\cdot\infty=\dfrac1\infty\cdot\infty$ 变成 $\dfrac\infty\infty$ 型, 而 $\dfrac\infty\infty=\dfrac10/\dfrac10=\dfrac00$, 等等. 下面我们利用柯西微分中值定理来给出这种不定式的定值方法 --- 洛必达 (L'Hospital) 法则. 以后我们还会发现, 虽然 $\dfrac\infty\infty$ 型可以化为 $\dfrac00$ 型, 但要利用洛必达法则, 我们还必须对 $\dfrac\infty\infty$ 的情况进行研究。

\subsection{$\dfrac00$ 型不定式}

\noindent\textbf{定理 5.2.1}\quad 设函数 $f(x)$ 和 $g(x)$ 在 $a$ 点的某一去心邻域 $U_0(a,\delta)$ 上可导, 而且满足:
\begin{enumerate}
  \item $\lim\limits_{x\to a}f(x)=\lim\limits_{x\to a}g(x)=0$;
  \item $g'(x)\ne0,\ \forall x\in U_0(a,\delta)$;
  \item $\lim\limits_{x\to a}\dfrac{f'(x)}{g'(x)}=\ell\ (\ell\text{ 为有限数或 }\pm\infty,\infty)$,
\end{enumerate}
则有
\[
  \lim_{x\to a}\frac{f(x)}{g(x)}
  =
  \lim_{x\to a}\frac{f'(x)}{g'(x)}
  =
  \ell.
\]

\noindent\textbf{证明}\quad 先考虑 $\ell$ 为有限数的情形, 并且我们先证 $\lim\limits_{x\to a-0}\dfrac{f(x)}{g(x)}=\ell$。
由条件 (1) 知, 若补充定义 $f(a)=g(a)=0$, 则 $f(x),g(x)$ 在 $a$ 点连续. 于是, 对 $\forall x\in(a-\delta,a)$, 在区间 $[x,a]$ 上应用柯西微分中值定理, 有
\[
  \frac{f(x)}{g(x)}
  =
  \frac{f(x)-f(a)}{g(x)-g(a)}
  =
  \frac{f'(\xi)}{g'(\xi)},\quad x<\xi<a.
\]
又由条件 (3), 可得
\[
  \lim_{x\to a-0}\frac{f'(\xi)}{g'(\xi)}
  =
  \lim_{x\to a-0}\frac{f'(x)}{g'(x)}
  =
  \ell.
\]
所以
\[
  \lim_{x\to a-0}\frac{f(x)}{g(x)}=\ell.
\]
同理可证 $\lim\limits_{x\to a+0}\dfrac{f(x)}{g(x)}=\ell$. 综合起来就有 $\lim\limits_{x\to a}\dfrac{f(x)}{g(x)}=\ell$。
当 $\ell=+\infty,-\infty$ 或 $\infty$ 时, 证明类似, 请读者作为练习自己给出。

\noindent\textbf{注}\quad 从定理的证明可以看出, 把 $x\to a$ 改为 $x\to a-0$ 或 $x\to a+0$, 上述定理的结论仍然成立。

当 $x\to\infty$ 时, 我们有

\noindent\textbf{定理 5.2.2}\quad 设函数 $f(x),g(x)$ 在 $U=\{x: |x|>a>0\}$ 上可导, 而且满足:
\begin{enumerate}
  \item $\lim\limits_{x\to\infty}f(x)=\lim\limits_{x\to\infty}g(x)=0$;
  \item $g'(x)\ne0,\ \forall x\in U$;
  \item $\lim\limits_{x\to\infty}\dfrac{f'(x)}{g'(x)}=\ell\ (\ell\text{ 为有限数或 }\pm\infty,\infty)$,
\end{enumerate}
则有
\[
  \lim_{x\to\infty}\frac{f(x)}{g(x)}=\ell.
\]

\noindent\textbf{证明}\quad 我们先证 $x\to+\infty$ 的情形. 作自变量变换 $x=\dfrac1t$, 则 $x\to+\infty$ 相应于 $t\to0+0$. 于是有
\[
  \lim_{t\to0+0}
  \frac{f'\left(\frac1t\right)}
       {g'\left(\frac1t\right)}
  =
  \lim_{x\to+\infty}\frac{f'(x)}{g'(x)},
\]
并且由条件 (1), 有
\[
  \lim_{t\to0+0}f\left(\frac1t\right)=0,\quad
  \lim_{t\to0+0}g\left(\frac1t\right)=0.
\]
应用定理 5.2.1 于开区间 $\left(0,\dfrac1a\right)$ 上新变量 $t$ 的函数 $f\left(\dfrac1t\right)$ 和 $g\left(\dfrac1t\right)$, 并注意到它们关于 $t$ 的导数为
\[
  f'\left(\frac1t\right)\left(-\frac1{t^2}\right),\quad
  g'\left(\frac1t\right)\left(-\frac1{t^2}\right),
\]
可得
\[
  \lim_{t\to0+0}
  \frac{f\left(\frac1t\right)}
       {g\left(\frac1t\right)}
  =
  \lim_{t\to0+0}
  \frac{f'\left(\frac1t\right)\left(-\frac1{t^2}\right)}
       {g'\left(\frac1t\right)\left(-\frac1{t^2}\right)}
  =
  \lim_{x\to+\infty}\frac{f'(x)}{g'(x)}
  =
  \ell.
\]
因此, 我们有
\[
  \lim_{x\to+\infty}\frac{f(x)}{g(x)}=\ell.
\]
同理可证 $x\to-\infty$ 的情形. 这样我们就证明了 $\lim\limits_{x\to\infty}\dfrac{f(x)}{g(x)}=\ell$ 成立。

\noindent\textbf{注}\quad 从定理的证明可以看出, 把 $x\to\infty$ 改为 $x\to-\infty$ 或 $x\to+\infty$, 上述定理的结论仍然成立。

\noindent\textbf{例 5.2.1}\quad 求极限
\[
  \lim_{x\to0}\frac{\sinh x-x}{\sin x-x}.
\]

\noindent\textbf{解}\quad 这是一个 $\dfrac00$ 型不定式. 应用洛必达法则, 可得
\[
  \lim_{x\to0}\frac{\sinh x-x}{\sin x-x}
  =
  \lim_{x\to0}\frac{\cosh x-1}{\cos x-1}
  =
  \lim_{x\to0}\frac{\sinh x}{-\sin x}
  =
  \lim_{x\to0}\frac{\cosh x}{-\cos x}
  =
  -1.
\]

\noindent\textbf{例 5.2.2}\quad 求极限
\[
  \lim_{x\to0+0}\frac{\sqrt{x}}{1-e^{2\sqrt{x}}}.
\]

\noindent\textbf{解}\quad 这也是一个 $\dfrac00$ 型不定式. 令 $t=\sqrt{x}$, 则有
\[
  \lim_{x\to0+0}\frac{\sqrt{x}}{1-e^{2\sqrt{x}}}
  =
  \lim_{t\to0+0}\frac{t}{1-e^{2t}}
  =
  \lim_{t\to0+0}\frac1{-2e^{2t}}
  =
  -\frac12.
\]

\noindent\textbf{例 5.2.3}\quad 求极限
\[
  \lim_{x\to+\infty}x\left(\frac\pi2-\arctan x\right).
\]

\noindent\textbf{解}\quad 这是一个 $0\cdot\infty$ 型不定式. 先将其转化为 $\dfrac00$ 型不定式, 然后应用洛必达法则, 可得
\[
  \lim_{x\to+\infty}x\left(\frac\pi2-\arctan x\right)
  =
  \lim_{x\to+\infty}
  \frac{\frac\pi2-\arctan x}{\frac1x}
  =
  \lim_{x\to+\infty}
  \frac{-\frac1{1+x^2}}{-\frac1{x^2}}
  =
  1.
\]

\subsection{$\dfrac\infty\infty$ 型不定式}

从理论上来讲, $\dfrac\infty\infty$ 型不定式可化为 $\dfrac00$ 型. 但是有时候这样做会使问题更为复杂, 致使问题的解决失败. 例如, 求 $\lim\limits_{x\to+\infty}\dfrac{\ln x}{x^\alpha}\ (\alpha>0)$. 如果将它化成
\[
  \frac{\ln x}{x^\alpha}
  =
  \frac{\frac1{x^\alpha}}{\frac1{\ln x}},
\]
则分子、分母分别求导后所得表达式比原来的函数还要复杂. 因此, 我们有必要专门研究求 $\dfrac\infty\infty$ 型不定式极限的方法。

\noindent\textbf{定理 5.2.3}\quad 设函数 $f(x),g(x)$ 在 $a$ 点的某一去心邻域 $U_0(a,\delta_0)\ (\delta_0>0)$ 上可导, 且满足:
\begin{enumerate}
  \item $\lim\limits_{x\to a}g(x)=\infty$;
  \item $g'(x)\ne0,\ \forall x\in U_0(a,\delta_0)$;
  \item $\lim\limits_{x\to a}\dfrac{f'(x)}{g'(x)}=\ell\ (\ell\text{ 有限数或 }\pm\infty,\infty)$,
\end{enumerate}
则有
\[
  \lim_{x\to a}\frac{f(x)}{g(x)}=\ell.
\]

\noindent\textbf{证明}\quad 只对 $\ell\in\mathbb{R}$ 和 $x\to a-0$ 的情形证明, 其他的情形请读者自己给出证明。

$\forall\varepsilon>0$, 由条件 (2) 和 (3) 知, $\exists\delta_1>0,0<\delta_1<\delta_0$, 当 $a-\delta_1<\xi<a$ 时, 有
\[
  \left|\frac{f'(\xi)}{g'(\xi)}-\ell\right|<\frac\varepsilon3.
\]
对已经取定的 $\delta_1$ 及 $x\in(a-\delta_1,a)$, 在 $[a-\delta_1,x]$ 上应用柯西微分中值定理, 存在 $\xi\in[a-\delta_1,x]$, 使
\[
  \frac{f(x)-f(x_1)}{g(x)-g(x_1)}-\ell
  =
  \frac{f'(\xi)}{g'(\xi)}-\ell,
\]
这里 $x_1=a-\delta_1,\ \xi\in(a-\delta_1,a)$. 上式可化为
\[
  f(x)-f(x_1)-\ell[g(x)-g(x_1)]
  =
  \left[\frac{f'(\xi)}{g'(\xi)}-\ell\right][g(x)-g(x_1)],
\]
整理得
\[
  f(x)-\ell g(x)
  =
  [f(x_1)-\ell g(x_1)]
  +
  \left[\frac{f'(\xi)}{g'(\xi)}-\ell\right][g(x)-g(x_1)].
\]
在上式两边同除以 $g(x)$, 可得
\[
  \frac{f(x)}{g(x)}-\ell
  =
  \left[\frac{f'(\xi)}{g'(\xi)}-\ell\right]\left[1-\frac{g(x_1)}{g(x)}\right]
  +
  \frac{f(x_1)-\ell g(x_1)}{g(x)}.
\]
又由于 $\lim\limits_{x\to a-0}g(x)=\infty$, 故对固定的 $x_1$, 有
\[
  \lim_{x\to a-0}\frac{f(x_1)-\ell g(x_1)}{g(x)}=0,\quad
  \lim_{x\to a-0}\frac{g(x_1)}{g(x)}=0.
\]
所以 $\exists\delta_2(0<\delta_2<\delta_1)$, 使得当 $a-\delta_2<x<a$ 时, 有
\[
  \left|\frac{f(x_1)-\ell g(x_1)}{g(x)}\right|<\frac\varepsilon2,\quad
  \left|\frac{g(x_1)}{g(x)}\right|<\frac12.
\]
令 $\delta=\min\{\delta_1,\delta_2\}$, 则当 $a-\delta<x<a$ 时, 有
\[
  \left|\frac{f(x)}{g(x)}-\ell\right|
  \le
  \left|\frac{f'(\xi)}{g'(\xi)}-\ell\right|
  \left|1-\frac{g(x_1)}{g(x)}\right|
  +
  \left|\frac{f(x_1)-\ell g(x_1)}{g(x)}\right|
  \le
  \frac32\cdot\frac\varepsilon3+\frac\varepsilon2
  =
  \varepsilon.
\]
因此 $\lim\limits_{x\to a-0}\dfrac{f(x)}{g(x)}=\ell$。

\noindent\textbf{注}\quad 把 $x\to a$ 改为 $x\to a-0$ 或 $x\to a+0$, 上述定理的结论仍然成立。

\noindent\textbf{定理 5.2.4}\quad 设函数 $f(x),g(x)$ 在 $U=\{x: |x|>h>0\}$ 上可导, 而且满足:
\[
  (1)\quad \lim_{x\to\infty}g(x)=\infty;
\]
\[
  (2)\quad g'(x)\ne0,\ \forall x\in U;
\]
\[
  (3)\quad \lim_{x\to\infty}\frac{f'(x)}{g'(x)}=\ell\ (\ell\text{ 有限数或 }\pm\infty,\infty),
\]
则有
\[
  \lim_{x\to\infty}\frac{f(x)}{g(x)}
  =
  \lim_{x\to\infty}\frac{f'(x)}{g'(x)}
  =
  \ell.
\]
本定理的证明类似于定理 5.2.2 的证明, 请读者作为练习自己补证。

\noindent\textbf{注}\quad 把 $x\to\infty$ 改为 $x\to-\infty$ 或 $x\to+\infty$, 上述定理的结论仍然成立。

\noindent\textbf{例 5.2.4}\quad 求极限
\[
  \lim_{x\to+\infty}\frac{\ln x}{x^\varepsilon},\quad \text{其中 }\varepsilon>0.
\]

\noindent\textbf{解}\quad 这是一个 $\dfrac\infty\infty$ 型不定式, 应用定理 5.2.4, 有
\[
  \lim_{x\to+\infty}\frac{\ln x}{x^\varepsilon}
  =
  \lim_{x\to+\infty}\frac{\frac1x}{\varepsilon x^{\varepsilon-1}}
  =
  \lim_{x\to+\infty}\frac1{\varepsilon x^\varepsilon}=0.
\]

\noindent\textbf{例 5.2.5}\quad 求极限
\[
  \lim_{x\to+\infty}\frac{x^\alpha}{e^x},\quad \text{其中 }\alpha>0.
\]

\noindent\textbf{解}\quad 这也是一个 $\dfrac\infty\infty$ 型不定式, 应用定理 5.2.4, 有
\[
  \lim_{x\to+\infty}\frac{x^\alpha}{e^x}
  =
  \lim_{x\to+\infty}\frac{\alpha x^{\alpha-1}}{e^x}
  =
  \cdots
  =
  \lim_{x\to+\infty}
  \frac{\alpha(\alpha-1)\cdots(\alpha-[\alpha])x^{\alpha-[\alpha]-1}}{e^x}
  =
  0.
\]

设 $I$ 是一个区间, $n$ 是一个非负整数, 在今后我们用记号 $C^nI$ 表示 $I$ 上所有具有 $n$ 阶连续导数的函数的集合; 用记号 $C^\infty I$ 表示 $I$ 上具有任意阶导数的函数的集合. 特别地, $C^0I$ 表示 $I$ 上所有连续函数的集合。

\noindent\textbf{例 5.2.6}\quad 设函数
\[
  f(x)=
  \begin{cases}
    e^{-\frac1{x^2}}, & x\ne0,\\
    0, & x=0.
  \end{cases}
\]
证明 $f(x)\in C^\infty(-\infty,+\infty)$。

\noindent\textbf{证明}\quad 当 $x\ne0$ 时,
\[
  f'(x)=\frac2{x^3}e^{-\frac1{x^2}}.
\]
当 $x=0$ 时,
\[
  f'(0)=\lim_{x\to0}\frac{f(x)-f(0)}x
  =
  \lim_{t\to\infty}\frac{t}{e^{t^2}}
  =
  0\quad \left(x=\frac1t\right),
\]
且
\[
  \lim_{x\to0}f'(x)
  =
  \lim_{t\to\infty}\frac{2t^3}{e^{t^2}}
  =
  0,
\]
所以 $f(x)\in C^1(-\infty,+\infty)$。

继续求导, 我们发现 $f^{(n)}(x)$ 具有如下形式:
\[
  f^{(n)}(x)=
  \begin{cases}
    P_{3n}\left(\dfrac1x\right)e^{-\frac1{x^2}}, & x\ne0,\\
    0, & x=0,
  \end{cases}
\]
其中 $P_{3n}(u)$ 为变元 $u$ 的 $3n$ 次多项式. 用数学归纳法可证上述断言为真. 事实上, 若它对 $n$ 成立, 则当 $x\ne0$ 时,
\[
  f^{(n+1)}(x)
  =
  \left[
    \frac2{x^3}P_{3n}\left(\frac1x\right)
    -
    \frac1{x^2}P'_{3n}\left(\frac1x\right)
  \right]e^{-\frac1{x^2}}
  =
  P_{3(n+1)}\left(\frac1x\right)e^{-\frac1{x^2}},
\]
而
\[
  f^{(n+1)}(0)
  =
  \lim_{x\to0}
  \frac{P_{3n}\left(\frac1x\right)e^{-\frac1{x^2}}}{x}
  =
  \lim_{t\to\infty}\frac{tP_{3n}(t)}{e^{t^2}}
  =
  0.
\]
由此推知 $f(x)$ 任意次可微, 即 $f(x)\in C^\infty(-\infty,+\infty)$。

\noindent\textbf{注}\quad 上例中的函数是一个重要的函数, 在今后的学习中还会多次遇到它. 这个函数可以用来说明很多问题。

\subsection{其他类型不定式}

通过变量替换, 我们总可把其他形式的不定式转化为 $\dfrac00$ 型和 $\dfrac\infty\infty$ 型, 具体采用什么样替换, 这要对具体问题做具体分析, 灵活对待。

\noindent\textbf{例 5.2.7}\quad 求极限
\[
  \lim_{x\to0+0}x^\alpha\ln x\quad(\alpha>0).
\]

\noindent\textbf{解}\quad 这是一个 $0\cdot\infty$ 型不定式. 我们有
\[
  \lim_{x\to0+0}x^\alpha\ln x
  =
  \lim_{x\to0+0}\frac{\ln x}{\frac1{x^\alpha}}
  =
  \lim_{x\to0+0}\frac{\frac1x}{-\frac{\alpha}{x^{\alpha+1}}}
  =
  0.
\]

\noindent\textbf{例 5.2.8}\quad 求极限
\[
  \lim_{x\to a}\left(\frac{a_1^x+a_2^x+\cdots+a_n^x}{n}\right)^{\frac1x}
  \quad (a_k>0,\ k=1,2,\cdots,n),
\]
这里 $a=0$ 或 $\pm\infty$。

\noindent\textbf{解}\quad 令
\[
  y=\left(\frac{a_1^x+a_2^x+\cdots+a_n^x}{n}\right)^{\frac1x},
  \quad
  \ln y=\frac1x\ln\left(\frac{a_1^x+a_2^x+\cdots+a_n^x}{n}\right).
\]
当 $x\to0$ 时, 应用洛必达法则得
\[
  \lim_{x\to0}\ln y
  =
  \frac1n\ln(a_1a_2\cdots a_n)
  =
  \ln\sqrt[n]{a_1a_2\cdots a_n},
\]
故 $\lim\limits_{x\to0}y=\sqrt[n]{a_1a_2\cdots a_n}$。

当 $x\to+\infty$ 时, 记 $M=\max\{a_1,a_2,\cdots,a_n\}$, 可得
\[
  \lim_{x\to+\infty}\ln y=\ln M,\quad
  \lim_{x\to+\infty}y=M.
\]
当 $x\to-\infty$ 时, 令 $m=\min\{a_1,a_2,\cdots,a_n\}$, 可得
\[
  \lim_{x\to-\infty}\ln y=\ln m,\quad
  \lim_{x\to-\infty}y=m.
\]

\noindent\textbf{例 5.2.9}\quad 设 $x_0\in(0,1)$, 定义 $x_n=\sin x_{n-1},\ n=1,2,\cdots$, 证明
\[
  \lim_{n\to\infty}\sqrt{n}x_n=\sqrt3.
\]

\noindent\textbf{证明}\quad 首先有
\[
  \lim_{x\to0}\left(\frac1{\sin^2x}-\frac1{x^2}\right)=\frac13.
\]
事实上,
\[
\begin{aligned}
  \lim_{x\to0}\left(\frac1{\sin^2x}-\frac1{x^2}\right)
  &=
  \lim_{x\to0}\frac{x^2-\sin^2x}{x^2\sin^2x}
  =
  \lim_{x\to0}\frac{x^2-\sin^2x}{x^4} \\
  &=
  \lim_{x\to0}\frac{2x-\sin2x}{4x^3}
  =
  \lim_{x\to0}\frac{1-\cos2x}{6x^2}
  =
  \frac13.
\end{aligned}
\]
又
\[
  0<x_{n+1}=\sin x_n<x_n,\quad n=1,2,\cdots,
\]
故 $x_n\to0$. 于是
\[
  \lim_{n\to\infty}\left(\frac1{x_{n+1}^2}-\frac1{x_n^2}\right)=\frac13,
\]
从而由第二章 §2.1 中例 2.1.13 的结果,
\[
  \lim_{n\to\infty}\frac1{nx_n^2}
  =
  \lim_{n\to\infty}\frac1n\sum_{k=1}^{n-1}
  \left(\frac1{x_{k+1}^2}-\frac1{x_k^2}\right)
  =
  \frac13.
\]
于是
\[
  \lim_{n\to\infty}\sqrt{n}x_n=\sqrt3.
\]

最后我们须指出的是:

(1) 并非所有的 $\dfrac00$ 或 $\dfrac\infty\infty$ 型不定式都可用洛必达法则求其极限. 例如, 容易证明
\[
  \lim_{x\to\infty}\frac{x+\sin x}{x}=1,
\]
但是
\[
  \lim_{x\to\infty}\frac{(x+\sin x)'}{(x)'}
  =
  \lim_{x\to\infty}(1+\cos x)
\]
不存在。

(2) 应用洛必达法则时, 每步必须验证 $\lim\dfrac{f'(x)}{g'(x)}$ 是否存在的条件. 例如, 显然有
\[
  \lim_{x\to\infty}\frac{x-\sin x}{x+\sin x}=1,
\]
但若盲目地使用洛必达法则, 就会得到
\[
  \lim_{x\to\infty}\frac{x-\sin x}{x+\sin x}
  =
  \lim_{x\to\infty}\frac{1-\cos x}{1+\cos x}
  =
  \lim_{x\to\infty}\frac{\sin x}{-\sin x}
  =
  -1
\]
的错误结论。

\section{泰勒公式}

用简单函数来逼近复杂函数, 这是一个重要的研究课题. 一般来说, 我们最熟悉而又十分简单的函数就是多项式函数. 在本节中, 我们将考虑两类多项式的逼近问题: 一类是在一点附近的逼近问题; 另一类则是在一个区间上的逼近问题。

设 $f(x)$ 是区间 $I$ 上的一个函数. 若我们用 $P_n(x)$ 去近似它, 则
\[
  R_n(x)=f(x)-P_n(x)
\]
就反映了近似的误差. 一个好的逼近就应该能够根据函数的性质求出所需多项式以及相应的控制误差, 并且还要使得误差尽可能地小。

\subsection{带佩亚诺余项的泰勒公式}

设函数 $f(x)$ 在 $U(x_0,\delta)$ 内有定义. 若 $f(x)$ 在点 $x_0$ 处连续, 则
\[
  f(x)-f(x_0)\to0\quad (x\to x_0),
\]
即
\[
  f(x)=f(x_0)+o(1)\quad (x\to x_0);
\]
若 $f(x)$ 在点 $x_0$ 处可导, 则
\[
  f(x)=f(x_0)+f'(x_0)(x-x_0)+o(x-x_0)\quad (x\to x_0).
\]
由此可知, 若用线性函数 $f(x_0)+f'(x_0)(x-x_0)$ 来近似 $f(x)$, 则在 $x_0$ 附近, 比用 $f(x_0)$ 来近似它误差要小得多. 若 $f(x)$ 在 $x_0$ 处具有高阶的导数, 则我们有

\noindent\textbf{定理 5.3.1}\quad 设函数 $f(x)$ 在 $x_0$ 处具有 $n(n\ge1)$ 阶导数, 则有
\[
  f(x)=f(x_0)+f'(x_0)(x-x_0)+\frac{f''(x_0)}{2!}(x-x_0)^2+\cdots
  +\frac{f^{(n)}(x_0)}{n!}(x-x_0)^n+o((x-x_0)^n)
  \quad (x\to x_0).
\]

\noindent\textbf{证明}\quad 要证上式成立, 只要证
\[
\begin{aligned}
&\lim_{x\to x_0}\frac1{(x-x_0)^n}
\left\{
f(x)-\left[
f(x_0)+f'(x_0)(x-x_0)+\frac{f''(x_0)}{2!}(x-x_0)^2 \right.\right.\\
&\left.\left.\hspace{8em}
\cdots+\frac{f^{(n)}(x_0)}{n!}(x-x_0)^n
\right]
\right\}=0
\end{aligned}
\]
即可. 对上式左边应用洛必达法则 $(n-1)$ 次, 可知它等于
\[
  \lim_{x\to x_0}\frac1{n!}
  \left[
    \frac{f^{(n-1)}(x)-f^{(n-1)}(x_0)}{x-x_0}
    -
    f^{(n)}(x_0)
  \right]
  =
  0.
\]
最后的等式成立, 是因为 $f^{(n)}(x_0)$ 存在. 定理证毕。

\noindent\textbf{注}\quad 当 $f(x)$ 在 $x_0$ 处具有 $n(n\ge1)$ 阶导数时, 我们把多项式
\[
  P_n(x)=f(x_0)+f'(x_0)(x-x_0)+\frac{f''(x_0)}{2!}(x-x_0)^2
  +\cdots+\frac{f^{(n)}(x_0)}{n!}(x-x_0)^n
\]
称为 $f(x)$ 在 $x_0$ 处的泰勒 (Taylor) 多项式, 而将
\[
  f(x)=f(x_0)+f'(x_0)(x-x_0)+\frac{f''(x_0)}{2!}(x-x_0)^2
  +\cdots+\frac{f^{(n)}(x_0)}{n!}(x-x_0)^n+R_n(x)
\]
称为 $f(x)$ 在 $x_0$ 处的泰勒公式, 其中 $R_n(x)=f(x)-P_n(x)$ 称为泰勒公式的余项. 特别地, 当 $x_0=0$ 时, 我们称此时的泰勒公式为麦克劳林 (Maclaurin) 公式。

定理 5.3.1 告诉我们, 当 $f(x)$ 在 $x_0$ 处存在 $n$ 阶导数时, 泰勒公式的余项为
\[
  R_n(x)=o((x-x_0)^n)\quad (x\to x_0).
\]
这种余项称为佩亚诺 (Peano) 余项。

值得指出的是, 若 $f(x)$ 在点 $x_0$ 存在 $n$ 阶导数, 且在点 $x_0$ 附近成立
\[
  f(x)=a_0+a_1(x-x_0)+\cdots+a_n(x-x_0)^n+o((x-x_0)^n)\quad (x\to x_0),
  \tag{5.3.1}
\]
则必有
\[
  a_k=\frac{f^{(k)}(x_0)}{k!}\quad (k=0,1,2,\cdots,n).
\]
事实上, 在 (5.3.1) 式中令 $x\to x_0$, 即得 $a_0=f(x_0)$. 现将 (5.3.1) 式改写并逐次除以 $(x-x_0),(x-x_0)^2,\cdots$, 再令 $x\to x_0$, 即可依次推出
\[
  a_1=f'(x_0),\quad a_2=\frac{f''(x_0)}{2!},\quad \cdots,\quad
  a_k=\frac{f^{(k)}(x_0)}{k!}.
\]
以上论述说明了 $f(x)$ 在 $x_0$ 处的多项式逼近的某种唯一性, 它将给我们寻找 $f(x)$ 的泰勒公式提供许多方便。

\noindent\textbf{例 5.3.1}\quad 求 $f(x)=e^x$ 的带佩亚诺余项的麦克劳林公式。

\noindent\textbf{解}\quad 由 $f^{(k)}(x)=e^x$, 且 $e^0=1$, 得
\[
  a_k=\frac{f^{(k)}(0)}{k!}=\frac1{k!},\quad k=0,1,\cdots,n,
\]
所以
\[
  e^x=1+\frac{x}{1!}+\frac{x^2}{2!}+\cdots+\frac{x^n}{n!}+o(x^n)\quad (x\to0).
\]

\noindent\textbf{例 5.3.2}\quad 求 $f(x)=\sin x$ 的带佩亚诺余项的麦克劳林公式。

\noindent\textbf{解}\quad 由 $f^{(k)}(x)=\sin\left(x+\dfrac{k\pi}{2}\right)$, 得
\[
  f^{(2k)}(0)=0,\quad f^{(2k+1)}(0)=(-1)^k,\quad k=0,1,\cdots,n,
\]
所以
\[
  \sin x=x-\frac{x^3}{3!}+\frac{x^5}{5!}+\cdots+(-1)^{n-1}\frac{x^{2n-1}}{(2n-1)!}+o(x^{2n})\quad (x\to0).
\]

\noindent\textbf{例 5.3.3}\quad 求 $f(x)=\cos x$ 的带佩亚诺余项的麦克劳林公式。

\noindent\textbf{解}\quad 由 $f^{(k)}(x)=\cos\left(x+\dfrac{k\pi}{2}\right)$, 得
\[
  f^{(2k)}(0)=(-1)^k,\quad f^{(2k+1)}(0)=0,\quad k=0,1,\cdots,n,
\]
所以
\[
  \cos x=1-\frac{x^2}{2!}+\frac{x^4}{4!}+\cdots+(-1)^n\frac{x^{2n}}{(2n)!}+o(x^{2n+1})\quad (x\to0).
\]

\noindent\textbf{例 5.3.4}\quad 求 $f(x)=\ln(1+x)$ 的带佩亚诺余项的麦克劳林公式。

\noindent\textbf{解}\quad 由
\[
  f^{(k)}(x)=(-1)^{k-1}\frac{(k-1)!}{(1+x)^k},
\]
得
\[
  f(0)=0,\quad f^{(k)}(0)=(-1)^{k-1}(k-1)!,\quad k=1,2,\cdots,n,
\]
所以
\[
  \ln(1+x)=x-\frac{x^2}{2}+\frac{x^3}{3}+\cdots+(-1)^{n-1}\frac{x^n}{n}+o(x^n)\quad (x\to0).
\]

\noindent\textbf{例 5.3.5}\quad 求 $f(x)=(1+x)^\alpha\ (\alpha\in\mathbb{R})$ 的带佩亚诺余项的麦克劳林公式。

\noindent\textbf{解}\quad 由
\[
  f^{(k)}(x)=\alpha(\alpha-1)\cdots(\alpha-k+1)(1+x)^{\alpha-k},
\]
得
\[
  a_0=1,\quad
  a_k=\frac{\alpha(\alpha-1)\cdots(\alpha-k+1)}{k!},\quad k=1,2,\cdots,n,
\]
所以
\[
  (1+x)^\alpha
  =
  1+\alpha x+\frac{\alpha(\alpha-1)}{2!}x^2+\cdots
  +\frac{\alpha(\alpha-1)\cdots(\alpha-n+1)}{n!}x^n+o(x^n)
  \quad (x\to0).
\]

例 5.3.1--例 5.3.5 中函数的麦克劳林公式在今后的各章节及后续课程中经常要用到, 读者应牢牢地记住它们。

\noindent\textbf{例 5.3.6}\quad 求 $f(x)=\sin^2(1+x^2)$ 的带佩亚诺余项的麦克劳林公式。

\noindent\textbf{解}\quad 由于
\[
\begin{aligned}
  \sin^2(1+x^2)
  &=
  \frac12-\frac12\cos(2+2x^2) \\
  &=
  \frac12-\frac12\bigl(\cos2\cos(2x^2)-\sin2\sin(2x^2)\bigr),
\end{aligned}
\]
而
\[
  \cos(2x^2)=1-\frac{(2x^2)^2}{2!}+\frac{(2x^2)^4}{4!}+\cdots+(-1)^n\frac{(2x^2)^{2n}}{(2n)!}+o(x^{4n+2}),
\]
\[
  \sin(2x^2)=2x^2-\frac{(2x^2)^3}{3!}+\cdots+(-1)^{n-1}\frac{(2x^2)^{2n-1}}{(2n-1)!}+o(x^{4n}),
\]
因此
\[
\begin{aligned}
  \sin^2(1+x^2)
  &=
  \frac12-\frac12\cos2+(\sin2)x^2+(\cos2)x^4-\frac{2^2\sin2}{3!}x^6
  -\frac{2^3\cos2}{4!}x^8 \\
  &\quad+\cdots+(-1)^{n-1}\frac{2^{2n-2}\sin2}{(2n-1)!}x^{4n-2}
  +(-1)^{n+1}\frac{2^{2n-1}\cos2}{(2n)!}x^{4n}
  +o(x^{4n})
  \quad (x\to0).
\end{aligned}
\]

由于 $f'(x)$ 的 $n$ 阶导数是 $f(x)$ 的 $n+1$ 阶导数, 因此若
\[
  f'(x)=a_1+a_2x+a_3x^2+\cdots+a_nx^{n-1}+o(x^{n-1})\quad (x\to0),
\]
则
\[
  f(x)=f(0)+a_1x+\frac{a_2}{2}x^2+\frac{a_3}{3}x^3+\cdots+\frac{a_n}{n}x^n+o(x^n)\quad (x\to0).
\]
下面我们利用这一结果, 导出 $\arctan x$ 和 $\arcsin x$ 的麦克劳林公式。

\noindent\textbf{例 5.3.7}\quad 求 $\arctan x$ 的带佩亚诺余项的麦克劳林公式。

\noindent\textbf{解}\quad 由于
\[
  (\arctan x)'=\frac1{1+x^2}=1-x^2+x^4+\cdots+(-1)^nx^{2n}+o(x^{2n+1})\quad (x\to0),
\]
因此
\[
  \arctan x=x-\frac13x^3+\frac15x^5+\cdots+\frac{(-1)^n}{2n+1}x^{2n+1}+o(x^{2n+2})\quad (x\to0).
\]

\noindent\textbf{例 5.3.8}\quad 求 $\arcsin x$ 的带佩亚诺余项的麦克劳林公式。

\noindent\textbf{解}\quad 由于
\[
  (\arcsin x)'=(1-x^2)^{-\frac12}
  =
  1+\frac1{2!!}x^2+\frac{3!!}{4!!}x^4+\cdots+
  \frac{(2n-1)!!}{(2n)!!}x^{2n}+o(x^{2n+1})
  \quad (x\to0),
\]
因此
\[
  \arcsin x
  =
  x+\frac13\frac1{2!!}x^3+\frac15\frac{3!!}{4!!}x^5+\cdots
  +\frac1{2n+1}\frac{(2n-1)!!}{(2n)!!}x^{2n+1}
  +o(x^{2n+2})
  \quad (x\to0).
\]

下面我们再举几个例子来说明如何求一个函数的泰勒公式。

\noindent\textbf{例 5.3.9}\quad 求 $e^{\cos x}$ 的带佩亚诺余项的麦克劳林公式 (展到 $x^4$)。

\noindent\textbf{解}\quad 利用 $e^u$ 在 $u=0$ 及 $\cos x$ 在 $x=0$ 的泰勒公式,
\[
\begin{aligned}
  e^{\cos x}
  &=
  e\cdot e^{\cos x-1} \\
  &=
  e\left[1+\cos x-1+\frac{(\cos x-1)^2}{2!}+o((\cos x-1)^2)\right] \\
  &=
  e\left[1+\left(-\frac{x^2}{2!}+\frac{x^4}{4!}+o(x^4)\right)
  +\frac1{2!}\left(-\frac{x^2}{2!}+o(x^2)\right)^2+o(x^4)\right] \\
  &=
  e-\frac e2x^2+\frac e6x^4+o(x^4)\quad (x\to0).
\end{aligned}
\]

\noindent\textbf{例 5.3.10}\quad 求 $e^x\ln(1+x)$ 的带佩亚诺余项的麦克劳林公式 (展到 $x^4$)。

\noindent\textbf{解}\quad 我们利用 $e^x$ 和 $\ln(1+x)$ 的麦克劳林公式得
\[
\begin{aligned}
  e^x\ln(1+x)
  &=
  \left(1+x+\frac{x^2}{2!}+\frac{x^3}{3!}+o(x^3)\right)
  \left(x-\frac{x^2}{2}+\frac{x^3}{3}-\frac{x^4}{4}+o(x^4)\right) \\
  &=
  x+\frac{x^2}{2}+\frac{x^3}{3}+o(x^4)\quad (x\to0).
\end{aligned}
\]

下面我们来给出泰勒公式的一些简单应用。

\noindent\textbf{例 5.3.11}\quad 求极限
\[
  \lim_{x\to+\infty}\left(\sqrt[3]{x^3+3x}-\sqrt{x^2-2x}\right).
\]

\noindent\textbf{解}\quad 这是一个 $\infty-\infty$ 型不定式. 下面我们用泰勒展开来求极限。
\[
\begin{aligned}
  &\lim_{x\to+\infty}\left(\sqrt[3]{x^3+3x}-\sqrt{x^2-2x}\right) \\
  &=
  \lim_{x\to+\infty}x
  \left\{
    \left[1+\frac3{x^2}\right]^{\frac13}
    -
    \left[1-\frac2x\right]^{\frac12}
  \right\} \\
  &=
  \lim_{x\to+\infty}x
  \left[
    \frac1{x^2}+\frac1x+o\left(\frac1x\right)
  \right]
  =
  1.
\end{aligned}
\]

\noindent\textbf{例 5.3.12}\quad 求极限
\[
  \lim_{n\to\infty}n^2\left(1-n\sin\frac1n\right).
\]

\noindent\textbf{解}\quad 将原式恒等变形, 利用泰勒公式得
\[
\begin{aligned}
  \lim_{n\to\infty}n^2\left(1-n\sin\frac1n\right)
  &=
  \lim_{n\to\infty}n^3\left(\frac1n-\sin\frac1n\right) \\
  &=
  \lim_{n\to\infty}n^3
  \left\{
    \frac1n-
    \left[\frac1n-\frac1{3!}\frac1{n^3}+o\left(\frac1{n^4}\right)\right]
  \right\}
  =
  \frac16.
\end{aligned}
\]

\subsection{带拉格朗日余项的泰勒公式}

上面我们讨论了函数在一点附近用多项式逼近的问题, 相应的误差只给出定性描述, 不能具体估计误差的大小, 所以它只适用于求无穷小量的阶或求极限等问题. 若要具体计算函数值并且达到预先指定的误差, 就需要给出误差的定量描述. 带拉格朗日余项的泰勒公式就是为了解决这一问题而产生的。

\noindent\textbf{定理 5.3.2}\quad 设 $f(x)\in C^n[a,b]$, 而且在 $(a,b)$ 内存在 $n+1$ 阶导数, 则对任意 $x,x_0\in[a,b]$, 有
\[
\begin{aligned}
  f(x)
  &=
  f(x_0)+\frac{f'(x_0)}{1!}(x-x_0)+\frac{f''(x_0)}{2!}(x-x_0)^2
  +\cdots+\frac{f^{(n)}(x_0)}{n!}(x-x_0)^n \\
  &\quad+
  \frac{f^{(n+1)}(\xi)}{(n+1)!}(x-x_0)^{n+1},
\end{aligned}
  \tag{5.3.2}
\]
其中 $\xi$ 介于 $x$ 与 $x_0$ 之间。

\noindent\textbf{证明}\quad 作两个辅助函数:
\[
  F(t)=f(x)-\left[
  f(t)+\frac{f'(t)}{1!}(x-t)+\frac{f''(t)}{2!}(x-t)^2
  +\cdots+\frac{f^{(n)}(t)}{n!}(x-t)^n
  \right]
\]
和
\[
  G(t)=(x-t)^{n+1}.
\]
容易验证它们在相应闭区间上连续, 在相应开区间内可导, 且 $F(x)=G(x)=0$. 应用柯西微分中值定理, 并且注意到
\[
  F'(t)=-\frac{f^{(n+1)}(t)}{n!}(x-t)^n,\quad
  G'(t)=-(n+1)(x-t)^n,
\]
便有
\[
  \frac{F(x_0)}{G(x_0)}
  =
  \frac{F'(\xi)}{G'(\xi)}
  =
  \frac{f^{(n+1)}(\xi)}{(n+1)!},
\]
这里 $\xi$ 介于 $x$ 与 $x_0$ 之间. 由此即得结论。

\noindent\textbf{注}\quad (5.3.2) 中的余项
\[
  \frac{f^{(n+1)}(\xi)}{(n+1)!}(x-x_0)^{n+1}
\]
称为拉格朗日余项. 另外, 我们可将上述定理中的 $\xi$ 表为
\[
  \xi=x_0+\theta(x-x_0),\quad \text{其中 }0<\theta<1.
\]

下面我们列出几个常见函数的带拉格朗日余项的泰勒公式:
\[
  e^x=1+x+\frac{x^2}{2!}+\cdots+\frac{x^n}{n!}+\frac{e^{\theta x}}{(n+1)!}x^{n+1}
  \quad (-\infty<x<+\infty,\ 0<\theta<1);
\]
\[
  \sin x=x-\frac{x^3}{3!}+\cdots+(-1)^{n-1}\frac{x^{2n-1}}{(2n-1)!}
  +(-1)^n\frac{\cos\theta x}{(2n+1)!}x^{2n+1}
  \quad (-\infty<x<+\infty,\ 0<\theta<1);
\]
\[
  \cos x=1-\frac{x^2}{2!}+\frac{x^4}{4!}+\cdots+(-1)^n\frac{x^{2n}}{(2n)!}
  +(-1)^{n+1}\frac{\cos\theta x}{(2n+2)!}x^{2n+2}
  \quad (-\infty<x<+\infty,\ 0<\theta<1);
\]
\[
  \ln(1+x)=x-\frac{x^2}{2}+\cdots+(-1)^{n-1}\frac{x^n}{n}
  +(-1)^n\frac{x^{n+1}}{(n+1)(1+\theta x)^{n+1}}
  \quad (|x|<1,\ 0<\theta<1);
\]
\[
\begin{aligned}
  (1+x)^\alpha
  &=
  1+\alpha x+\frac{\alpha(\alpha-1)}{2!}x^2+\cdots
  +\frac{\alpha(\alpha-1)\cdots(\alpha-n+1)}{n!}x^n \\
  &\quad+
  \frac{\alpha(\alpha-1)\cdots(\alpha-n)}{(n+1)!}(1+\theta x)^{\alpha-n-1}x^{n+1}
  \quad (|x|<1,\ 0<\theta<1).
\end{aligned}
\]
利用拉格朗日余项, 我们可以估计在整个区间上 $f(x)$ 用泰勒公式逼近时的误差. 这在实际应用中是十分重要的。

\noindent\textbf{例 5.3.13}\quad 求 $e$ 的近似值, 使得其误差不超过 $10^{-5}$。

\noindent\textbf{解}\quad 在 $e^x$ 的带拉格朗日余项的泰勒公式中令 $x=1$, 得
\[
  e=1+1+\frac1{2!}+\cdots+\frac1{n!}+\frac{e^\theta}{(n+1)!},\quad 0<\theta<1.
\]
欲使
\[
  \frac{e^\theta}{(n+1)!}<\frac3{(n+1)!}<10^{-5},
\]
只需取 $n=8$ 即可. 于是, 有
\[
  e\approx1+1+\frac1{2!}+\cdots+\frac1{8!}\approx2.71828.
\]

\noindent\textbf{例 5.3.14}\quad 试证 $e$ 是无理数。

\noindent\textbf{证明}\quad 用反证法. 假设 $e$ 是有理数, 它表示成分数时的分母为 $N$. 现取正整数 $n>\max\{N,3\}$. 根据 $e^x$ 的带拉格朗日余项的泰勒公式, 得
\[
  e-\left(1+1+\frac1{2!}+\cdots+\frac1{n!}\right)=\frac{e^\theta}{(n+1)!}.
\]
这里 $0<\theta<1$, 从而 $1<e^\theta<3$. 上式两边乘以 $n!$, 得
\[
  n!\left[e-\left(1+1+\frac1{2!}+\cdots+\frac1{n!}\right)\right]
  =
  \frac{e^\theta}{n+1}.
\]
上式左边是整数, 而右边满足 $0<\dfrac{e^\theta}{n+1}<1$. 这一矛盾说明 $e$ 是无理数。

\noindent\textbf{例 5.3.15}\quad 求极限
\[
  \lim_{n\to\infty}n\sin(2\pi en!).
\]

\noindent\textbf{解}\quad 根据 $e^x$ 的带拉格朗日余项的泰勒公式, 得
\[
  n!e
  =
  n!\left(1+1+\frac1{2!}+\cdots+\frac1{n!}+\frac1{(n+1)!}+\frac{e^\theta}{(n+2)!}\right)
  =
  k+\frac1{n+1}+\frac{e^\theta}{(n+2)(n+1)},
\]
其中 $0<\theta<1$, 而
\[
  k=n!\left(1+1+\frac1{2!}+\cdots+\frac1{n!}\right)
\]
是正整数. 于是
\[
\begin{aligned}
  \lim_{n\to\infty}n\sin(2\pi en!)
  &=
  \lim_{n\to\infty}
  n\sin\left(2k\pi+\frac{2\pi}{n+1}
  +\frac{2\pi e^\theta}{(n+2)(n+1)}\right) \\
  &=
  \lim_{n\to\infty}
  n\left[\frac{2\pi}{n+1}+o\left(\frac1n\right)\right]
  =
  2\pi.
\end{aligned}
\]

\subsection{拉格朗日插值多项式}

在实际应用中大量的函数都是用表格给出的, 为了进行理论研究和工程设计, 需要寻找与已知函数值相符而形式简单的函数去近似它. 解决这类问题的最基本方法之一就是多项式插值法. 设 $y=f(x)$ 是在区间 $[a,b]$ 上定义的某个函数, 已知它在该区间上 $n+1$ 个不同点 $x_0,x_1,x_2,\cdots,x_n$ 处的函数值为
\[
  y_i=f(x_i)\quad (i=0,1,2,\cdots,n),
\]
寻找一个多项式 $P(x)$, 使得
\[
  P(x_i)=y_i\quad (i=0,1,2,\cdots,n).
  \tag{5.3.2}
\]
对于给定的插值数据 $\{(x_i,y_i)\}_{i=0}^n$, 满足条件 (5.3.2) 的多项式 $P(x)$ 称为 $f(x)$ 的插值多项式。

先来看 $n=1$ 的情形. 通过 $(x_0,y_0)$ 和 $(x_1,y_1)$ 两点的直线可表示为
\[
  L_1(x)=y_0\cdot\frac{x-x_1}{x_0-x_1}+y_1\cdot\frac{x-x_0}{x_1-x_0}.
\]
其中
\[
  \ell_0(x)=\frac{x-x_1}{x_0-x_1},\quad
  \ell_1(x)=\frac{x-x_0}{x_1-x_0},
\]
满足
\[
  \ell_0(x_0)=1,\quad \ell_0(x_1)=0,\quad \ell_1(x_0)=0,\quad \ell_1(x_1)=1.
\]
若能构造出具有如下性质的特殊插值多项式 $\ell_i(x)$:
\[
  \ell_i(x_j)=\delta_{ij}=
  \begin{cases}
    1, & i=j,\\
    0, & i\ne j,
  \end{cases}
\]
则多项式
\[
  P(x)=\sum_{i=0}^n y_i\cdot\ell_i(x)
\]
就是满足插值条件的插值多项式. 注意到 $x_0,\cdots,x_n$ 中除 $x_i$ 外, 均为多项式 $\ell_i(x)$ 的零点, 即知
\[
  \ell_i(x)
  =
  \frac{(x-x_0)\cdots(x-x_{i-1})(x-x_{i+1})\cdots(x-x_n)}
       {(x_i-x_0)\cdots(x_i-x_{i-1})(x_i-x_{i+1})\cdots(x_i-x_n)}.
\]
我们把上面所确定的插值多项式记为
\[
  L_n(x)=\sum_{i=0}^n y_i\frac{\omega(x)}{(x-x_i)\omega'(x_i)},
\]
其中
\[
  \omega(x)=(x-x_0)(x-x_1)\cdots(x-x_n),
  \tag{5.3.3}
\]
并且称之为 $n$ 阶拉格朗日插值多项式, 而称 $\{\ell_i(x)\}_{i=0}^n$ 为拉格朗日插值基函数. 另外, 我们称 $x_0,x_1,\cdots,x_n$ 为插值节点。

\noindent\textbf{例 5.3.16}\quad 已知函数 $f(x)$ 在三点的函数值为
\[
\begin{array}{c|ccc}
  x & 1.0 & 1.5 & 2.0\\
  \hline
  f(x) & 0.0000 & 0.4055 & 0.6931
\end{array}
\]
求它的二阶插值多项式 $P_2(x)$, 并用 $P_2(x)$ 估算 $f(1.2)$。

\noindent\textbf{解}\quad 记 $x_0=1.0,x_1=1.5,x_2=2.0$, 则有
\[
\begin{aligned}
  L_2(x)
  &=
  0.0000\times\frac{(x-1.5)(x-2)}{(1-1.5)(1-2)}
  +0.4055\times\frac{(x-1)(x-2)}{(1.5-1)(1.5-2)} \\
  &\quad
  +0.6931\times\frac{(x-1)(x-1.5)}{(2-1)(2-1.5)} \\
  &=
  -1.6220(x-1)(x-2)+1.3862(x-1)(x-1.5).
\end{aligned}
\]
取 $x=1.2$, 得
\[
\begin{aligned}
  f(1.2)\approx L_2(1.2)
  &=
  -1.6220(1.2-1)(1.2-2)
  +1.3862(1.2-1)(1.2-1.5) \\
  &=
  0.176348.
\end{aligned}
\]
由于原始数据只有 $4$ 位小数, 我们可以取 $f(1.2)=0.1763$。

用插值多项式 $L_n(x)$ 作为被插函数 $f(x)$ 的逼近, 除了在插值节点处以外, $L_n(x)$ 与被插函数 $f(x)$ 是有差别的, 下面的定理给出了插值余项 (误差) 的一个表达式。

\noindent\textbf{定理 5.3.3}\quad 设被插函数 $f(x)\in C^{(n+1)}[a,b]$, 且插值节点互不相同, 则对任意的 $x\in[a,b]$, 都存在 $\xi\in[a,b]$, 使得
\[
  R_n(x)=f(x)-L_n(x)=\frac{f^{(n+1)}(\xi)}{(n+1)!}\omega(x),
\]
其中 $\omega(x)$ 如 (5.3.3) 所定义。

\noindent\textbf{证明}\quad 由于
\[
  f(x_j)=L_n(x_j),\quad j=0,1,2,\cdots,n,
\]
所以 $x_0,x_1,\cdots,x_n$ 为误差函数 $R_n(x)$ 的零点. 于是
\[
  R_n(x)=K(x)\omega(x).
\]
为给出 $K(x)$ 的表达式, 引入变量为 $t$ 的函数
\[
  F(t)=f(t)-L_n(t)-K(x)\omega(t).
\]
任取 $x\in[a,b]$, 且 $x\ne x_j$, 则有 $F(x)=0,\ F(x_j)=0$. 反复应用罗尔微分中值定理, 可得 $F^{(n+1)}(t)$ 至少有一个零点 $\xi$, 即
\[
  0=F^{(n+1)}(\xi)=f^{(n+1)}(\xi)-K(x)(n+1)!.
\]
因此
\[
  K(x)=\frac{f^{(n+1)}(\xi)}{(n+1)!}.
\]

\section{利用导数研究函数}

这一节我们将利用导数来研究函数的单调性、极值、凹凸性和拐点等. 最后再利用函数的这些特性, 给出描绘函数图像的方法。

\subsection{函数的单调性}

这一小节, 作为拉格朗日中值微分定理的应用, 我们来讨论如何利用导数的符号来判断函数的单调性. 事实上, 从几何上来看, 函数的升降, 与切线的方向密切相关. 曲线单调上升时, 切线与 $x$ 轴正向的夹角为锐角; 曲线单调下降时, 切线与 $x$ 轴正向的夹角为钝角. 这样, 我们便有

\noindent\textbf{定理 5.4.1}\quad 设函数 $f(x)$ 在区间 $I$ 内可导, 则

(1) $f(x)$ 在 $I$ 内单调上升的充分必要条件是 $f'(x)\ge0,\ \forall x\in I$;

(2) $f(x)$ 在 $I$ 内单调下降的充分必要条件是 $f'(x)\le0,\ \forall x\in I$。

\noindent\textbf{证明}\quad 我们只证 (1). (2) 的证明是类似的。

当 $f(x)$ 在 $I$ 内单调上升时, 对任意的 $x\in I$, 取充分小的 $\Delta x>0$ 使得 $x+\Delta x\in I$. 由于 $f(x+\Delta x)-f(x)\ge0$, 从而有
\[
  \frac{f(x+\Delta x)-f(x)}{\Delta x}\ge0.
\]
由于 $f(x)$ 在 $I$ 内可导, 因此有
\[
  f'(x)=f'(x+0)=\lim_{\Delta x\to0+0}\frac{f(x+\Delta x)-f(x)}{\Delta x}\ge0.
\]
这就证明了必要性。

现证充分性. 任取 $x_1,x_2\in I$, 并且 $x_1<x_2$, 在 $[x_1,x_2]$ 上应用拉格朗日微分中值定理, 知 $\exists\xi\in(x_1,x_2)$, 使得
\[
  f(x_2)-f(x_1)=f'(\xi)(x_2-x_1)\ge0.
\]
因此 $f(x_1)\le f(x_2)$, 即 $f(x)$ 在 $I$ 内单调上升。

\noindent\textbf{注}\quad 大家已经知道, 一个可微函数在一个区间上为常数的充分必要条件是该函数的导数在该区间上恒等于零. 因此我们立即可得: 一个可导函数 $f(x)$ 在 $I$ 内严格单调上升 (下降) 的充分必要条件是在区间 $I$ 上 $f'(x)\ge0\ (f'(x)\le0)$, 而且在该区间的任何子区间上 $f'(x)$ 都不恒等于 $0$。

\noindent\textbf{例 5.4.1}\quad 设函数
\[
  f(x)=
  \begin{cases}
    e^{\sin\frac1x-\frac1x}, & x>0,\\
    0, & x=0.
  \end{cases}
\]
对于任意的 $x>0$, 有
\[
  f'(x)=e^{\sin\frac1x-\frac1x}\left(1-\cos\frac1x\right)\frac1{x^2}\ge0,
\]
等号成立当且仅当 $x=\dfrac1{2k\pi}\ (k=1,2,\cdots)$. 此外, 由于
\[
  \lim_{x\to0+0}f(x)=0=f(0),
\]
所以 $f(x)$ 在 $[0,+\infty)$ 上连续. 这样, 应用定理 5.4.1 的注即知, $f(x)$ 在 $[0,+\infty)$ 上严格单调上升。

\noindent\textbf{例 5.4.2}\quad 证明不等式
\[
  \tan x>x+\frac13x^3\quad \left(0<x<\frac\pi2\right).
\]

\noindent\textbf{证明}\quad 令
\[
  f(x)=\tan x-x-\frac13x^3,
\]
则有
\[
  f'(x)=\frac1{\cos^2x}-1-x^2=\tan^2x-x^2>0,\quad \forall x\in\left(0,\frac\pi2\right),
\]
这是因为在该区间上 $\tan x>x$ 之故. 因此, $f(x)$ 在 $[0,\pi/2)$ 上严格单调上升. 又 $f(0)=0$, 故有
\[
  0=f(0)<f(x)=\tan x-x-\frac13x^3,\quad \forall x\in\left(0,\frac\pi2\right),
\]
从而不等式得证。

\subsection{函数的极值}

在本章的第一节我们已经介绍了极值的定义和费马定理. 我们已经知道, 一个在区间 $I$ 内可导的函数 $f(x)$, 若它在 $I$ 内的一点 $x_0$ 上取极值, 则必有 $f'(x_0)=0$ (即 $x_0$ 是 $f(x)$ 的一个驻点), 但反之不真. 这样, 假如我们已经知道 $x_0$ 是 $f(x)$ 的一个驻点, 那么在什么条件下它是 $f(x)$ 的一个极值点呢? 下面我们利用导数的符号来回答这一问题。

\noindent\textbf{定理 5.4.2}\quad 设函数 $f(x)$ 在 $U_0(x_0,\delta)$ 内可导且在点 $x_0$ 处连续。

(1) 若当 $x_0-\delta<x<x_0$ 时, $f'(x)>0$, 而当 $x_0<x<x_0+\delta$ 时, $f'(x)<0$, 则 $f(x)$ 在点 $x_0$ 取得严格极大值;

(2) 若当 $x_0-\delta<x<x_0$ 时, $f'(x)<0$, 而当 $x_0<x<x_0+\delta$ 时, $f'(x)>0$, 则 $f(x)$ 在点 $x_0$ 取得严格极小值;

(3) 若当 $x_0-\delta<x<x_0$ 及 $x_0<x<x_0+\delta$ 时, 都有 $f'(x)>0$, 或者 $f'(x)<0$, 则 $x_0$ 不是 $f(x)$ 的极值点。

请读者作为练习给出定理的证明。

定理 5.4.2 就是说, 当 $x$ 从点 $x_0$ 的左侧经过 $x_0$ 而变到其右侧时, 若 $f'(x)$ 改变符号, 则 $f(x)$ 在点 $x_0$ 取得极值: 由正变负, 取极大值; 由负变正, 取极小值. 若 $f'(x)$ 不改变符号, 则 $f(x)$ 在点 $x_0$ 不取极值. 对于上述定理中的条件与结论, 读者应该根据函数图像的几何特征及导数的几何意义来加以理解与记忆, 而不是简单地记住该定理. 如图 5.4.1 所示, $x_0$ 是 $f(x)$ 的极大值点, $x_1$ 是 $f(x)$ 的极小值点, $x_2$ 不是 $f(x)$ 的极值点。

\noindent\textbf{例 5.4.3}\quad 求函数 $f(x)=(x+2)^2(x-1)^3$ 在 $(-\infty,+\infty)$ 上的极值。

\noindent\textbf{解}\quad 直接计算, 得
\[
  f'(x)=(x+2)(x-1)^2(5x+4).
\]
由此可知, 该函数有 $3$ 个驻点:
\[
  x_1=-2,\qquad x_2=-\frac45,\qquad x_3=1.
\]
这 $3$ 个点将实数轴分成 $4$ 个区间:
\[
  (-\infty,-2),\qquad \left(-2,-\frac45\right),\qquad
  \left(-\frac45,1\right),\qquad (1,+\infty).
\]
容易看出, 导数 $f'(x)$ 在这 $4$ 个区间上的符号分别为:
\[
  +,\quad -,\quad +,\quad +.
\]
因此, 函数 $f(x)$ 在点 $x_1=-2$ 取得极大值 $0$; 在点 $x_2=-\dfrac45$ 取得极小值
\[
  f\left(-\frac45\right)=\left(\frac65\right)^2\left(-\frac95\right)^3\approx -8.4;
\]
而在点 $x_3=1$ 不取极值。

我们也可以根据函数在一点的高阶导数的符号来判定极值点. 我们有

\noindent\textbf{定理 5.4.3}\quad 设函数 $f(x)$ 在 $U(x_0,\delta)$ 内 $n$ 阶可导, $f'(x_0)=f''(x_0)=\cdots=f^{(n-1)}(x_0)=0$, 且 $f^{(n)}(x_0)\ne0$, 则

(1) 当 $n$ 为奇数时, $f(x)$ 在点 $x_0$ 不取极值;

(2) 当 $n$ 为偶数且 $f^{(n)}(x_0)>0$ 时, $f(x)$ 在 $x_0$ 取严格极小值;

(3) 当 $n$ 为偶数且 $f^{(n)}(x_0)<0$ 时, $f(x)$ 在 $x_0$ 取严格极大值。

\noindent\textbf{证明}\quad 由于函数 $f(x)$ 在 $U(x_0,\delta)$ 内 $n$ 阶可导, 而且
\[
  f'(x_0)=f''(x_0)=\cdots=f^{(n-1)}(x_0)=0,
\]
利用带佩亚诺余项的泰勒公式, 可得
\begin{equation}
\begin{aligned}
  f(x)-f(x_0)
  &=\frac{f^{(n)}(x_0)}{n!}(x-x_0)^n+o((x-x_0)^n)\\
  &=\left(\frac{f^{(n)}(x_0)}{n!}+o(1)\right)(x-x_0)^n\quad (x\to x_0).
\end{aligned}
\tag{5.4.1}
\end{equation}
而
\[
  \lim_{x\to x_0}\left(\frac{f^{(n)}(x_0)}{n!}+o(1)\right)
  =\frac{f^{(n)}(x_0)}{n!},
\]
因此存在 $\delta_1\in(0,\delta)$, 使得当 $x\in U(x_0,\delta_1)$ 时, 有 $\dfrac{f^{(n)}(x_0)}{n!}+o(1)$ 与 $f^{(n)}(x_0)$ 同号. 于是, 我们有:

(1) 当 $n$ 为奇数时, 若 $x$ 从点 $x_0$ 的左侧经过 $x_0$ 而变到其右侧, 等式 (5.4.1) 的右边将改变符号, 从而 $f(x)-f(x_0)$ 也改变符号. 这表明在 $x_0$ 的任何邻域内总有两点 $x_1$ 和 $x_2$, 使得
\[
  (f(x_0)-f(x_1))(f(x_0)-f(x_2))<0.
\]
因此 $f(x)$ 在 $x_0$ 不取极值;

(2) 当 $n$ 为偶数且 $f^{(n)}(x_0)>0$ 时, 对 $\forall x\in U_0(x_0,\delta_1)$, (5.4.1) 式的右边值大于零, 从而有 $f(x)>f(x_0)$, 即 $f(x)$ 在 $x_0$ 取严格极小值; 而当 $n$ 为偶数且 $f^{(n)}(x_0)<0$ 时, $\forall x\in U_0(x_0,\delta_1)$, (5.4.1) 式的右边值小于零, 从而有 $f(x)<f(x_0)$, 即 $f(x)$ 在 $x_0$ 取严格极大值. 定理证毕。

\noindent\textbf{注}\quad 对于定理 5.4.3, 最常见的是 $n=2$ 的情形, 即若 $f'(x_0)=0$, 但 $f''(x_0)\ne0$, 则

(1) 当 $f''(x_0)>0$ 时, $f(x)$ 在 $x_0$ 取严格极小值;

(2) 当 $f''(x_0)<0$ 时, $f(x)$ 在 $x_0$ 取严格极大值。

若对函数 $y=x^n\ (n\in\mathbb{N})$ 来讨论其在点 $x=0$ 处的极值情况, 则定理 5.4.3 的结论是显然的. 事实上, 定理 5.4.3 告诉我们, 利用泰勒公式, 任何满足该定理条件的函数 $f(x)$ 在点 $x=x_0$ 处取极值情况和函数 $f^{(n)}(x_0)(x-x_0)^n$ 完全一样. 换句话说, 读者可利用 $ax^n(a\ne0)$ 在 $x=0$ 的取值情况来记忆定理 5.4.3。

如果我们利用定理 5.4.3 来解例 5.4.3, 则由 $f''(-2)=-54<0$ 立刻可以得出 $f(x)$ 在 $x_1=-2$ 取极大值; 由 $f''(1)=0$ 及 $f'''(1)=54>0$ 知 $f(x)$ 在 $x=1$ 不取极值. 再注意到 $f''\left(-\dfrac45\right)>0$ 即可推出 $f(x)$ 在 $x_2=-\dfrac45$ 也取极小值。

\noindent\textbf{例 5.4.4 (达布 (Darboux) 定理)}\quad 设函数 $f(x)$ 在 $[a,b]$ 上可导, 则对于任意介于 $f'(a+0)$ 与 $f'(b-0)$ 之间的数 $\eta$, 都存在 $\xi\in(a,b)$, 使得 $f'(\xi)=\eta$。

\noindent\textbf{证明}\quad 先假设 $f'(a+0)>0,\ f'(b-0)<0,\ \eta=0$. 根据导数的定义, 由 $f'(a+0)>0$ 可导出, 存在 $\delta_1>0$, 使得 $\forall x\in[a,a+\delta_1]$, 有
\[
  f(x)>f(a);
\]
再由 $f'(b-0)<0$ 可导出, 存在 $\delta_2>0$, 使得 $\forall x\in[b-\delta_2,b]$, 有
\[
  f(x)>f(b).
\]
此外, 由于 $f(x)$ 在 $[a,b]$ 上可导, 从而 $f(x)$ 在 $[a,b]$ 上连续. 现设
\[
  M=\max_{x\in[a,b]}\{f(x)\},
\]
则存在 $\xi\in[a,b]$, 使得 $f(\xi)=M$. 由以上分析可知 $\xi\ne a,b$, 因此费马定理告诉我们 $f'(\xi)=0$。

同理可证 $f'(a+0)<0,\ f'(b-0)>0$ 和 $\eta=0$ 的情形。

对一般情形, 我们构造辅助函数 $F(x)=f(x)-\eta x$. 由 $F'(a+0)=f'(a+0)-\eta$ 和 $F'(b-0)=f'(b-0)-\eta$ 异号知, 存在 $\xi\in(a,b)$, 使得
\[
  0=F'(\xi)=f'(\xi)-\eta.
\]

\noindent\textbf{注}\quad 值得注意的是, 一个函数即使在一个区间上可导, 其导函数也不一定是连续的 (参见下例). 而上述例子告诉我们, 即使导函数不连续, 仍然具有类似于连续函数的介值定理成立。

\noindent\textbf{例 5.4.5}\quad 设函数
\[
  f(x)=
  \begin{cases}
    x^2\sin\dfrac1x, & x\ne0,\\
    0, & x=0,
  \end{cases}
\]
则
\[
  f'(x)=
  \begin{cases}
    2x\sin\dfrac1x-\cos\dfrac1x, & x\ne0,\\
    0, & x=0.
  \end{cases}
\]
显然, $f'(x)$ 在 $x=0$ 处不连续。

\noindent\textbf{思考题}\quad 设函数 $f(x)$ 在区间 $I$ 上可导, 且 $f'(x)\ne0,\ \forall x\in I$. 试证明 $f(x)$ 是 $I$ 上的严格单调函数。

此外, 我们需要特别指出的是:

(1) 函数的极值只是函数的局部特性, 一个函数的某些极小值完全可以远远大于它的某些极大值;

(2) 函数的极值点必须位于函数所定义区间的内部, 区间的端点不涉及函数的极值问题。

在实际应用中我们常常遇到的是求一个函数在指定范围内的最大值或者最小值问题. 大家已经知道, 若函数 $f(x)$ 在闭区间 $[a,b]$ 上连续, 则 $f(x)$ 在 $[a,b]$ 上取得它的最大值和最小值. 现在我们可以通过比较函数 $f(x)$ 在 $[a,b]$ 上所有的驻点、不可导点和端点的函数值的大小来确定它的最大值和最小值。

\noindent\textbf{例 5.4.6}\quad 求函数 $f(x)=\sin^3x+\cos^3x$ 在 $\left[-\dfrac\pi4,\dfrac{3\pi}4\right]$ 上的最大值和最小值。

\noindent\textbf{解}\quad 直接计算, 得
\[
  f'(x)=3\sin x\cos x(\sin x-\cos x).
\]
于是, 在所考虑区间内有函数的 $3$ 个驻点:
\[
  x_1=0,\qquad x_2=\frac\pi4,\qquad x_3=\frac\pi2.
\]
易知
\[
  f\left(\frac\pi2\right)=f(0)=1,\quad
  f\left(\frac\pi4\right)=\frac{\sqrt2}{2},\quad
  f\left(-\frac\pi4\right)=f\left(\frac{3\pi}4\right)=0.
\]
因此, 该函数的最大值和最小值分别为 $1$ 和 $0$。

\noindent\textbf{例 5.4.7}\quad 作一个无盖圆柱形茶缸, 使得它的底的厚度是侧面厚度的三倍. 试问: 容积一定时如何确定它的底面半径和高才能使得用料最省?

\noindent\textbf{解}\quad 设茶缸的侧面的厚度为 $\delta$, 则它的底的厚度为 $3\delta$. 记茶缸的容积为 $V=$ 常数, 底面半径为 $r$, 则高为 $h=\dfrac{V}{\pi r^2}$. 于是, 若茶缸的用料量用 $\rho S(r)$ 表示, 其中 $\rho$ 为所用材料的密度, 则有
\[
  S(r)=\delta(3\pi r^2+2\pi rh)
  =\delta\left(3\pi r^2+\frac{2V}{r}\right).
\]
这样, 所考虑的问题就等价于求函数 $S(r)$ 的极小值。

对 $S(r)$ 求导, 得
\[
  S'(r)=\delta\left(6\pi r-\frac{2V}{r^2}\right).
\]
因此 $S'(r)$ 在 $[0,+\infty)$ 内有唯一的零点 $r_0=\sqrt[3]{\dfrac{V}{3\pi}}$. 又
\[
  S''(r)=\delta\left(6\pi+\frac{4V}{r^3}\right)>0,\quad \forall r\in(0,+\infty),
\]
所以 $r_0$ 是 $S(r)$ 的唯一极小值点. 这时, 对应的高为
\[
  h_0=\frac{V}{\pi r_0^2}
  =\frac{3\pi r_0^3}{\pi r_0^2}=3r_0.
\]
这也就是说, 当高为底面半径的 $3$ 倍时用料最省。

\subsection{函数的凹凸性}

函数的凹凸性是函数一个重要的几何特征. 设函数 $f(x)$ 在区间 $I$ 内有定义. 从几何上来看, 若 $y=f(x)$ 的图像上任意两点 $(x_1,f(x_1))$ 和 $(x_2,f(x_2))$ 之间的曲线段总位于连接这两点的线段之下 (上), 则称该函数是凸 (凹) 的 (参见图 5.4.2). 这个概念用解析的语言可以表述成定义 5.4.1。

\noindent\textbf{定义 5.4.1}\quad 设函数 $f(x)$ 在区间 $I$ 内有定义. 若 $\forall x_1,x_2\in I,\ \forall t\in(0,1)$, 有
\[
  f(tx_1+(1-t)x_2)\le tf(x_1)+(1-t)f(x_2),
\]
则称 $f(x)$ 为 $I$ 上的凸函数; 若当 $x_1\ne x_2$ 时, 总有严格不等式成立, 则称 $f(x)$ 为 $I$ 上的严格凸函数。

\noindent\textbf{注}\quad 在上述定义中将 ``$\le(<)$'' 改为 ``$\ge(>)$'', 就得到了凹 (严格凹) 函数的定义. 所以有一个凸函数的定理, 就有一个相对应的凹函数的定理. 下面我们只讨论凸函数。

设 $f(x)$ 在 $\mathbb{R}$ 中有定义, 容易看出 $f(x)$ 既是凸函数又是凹函数的充分必要条件是 $f(x)$ 是一个线性函数; 若 $f(x)=ax^2+bx+c\ (a,b,c\in\mathbb{R})$, 则当 $a>0$ 时, $f(x)$ 是凸函数; 而当 $a<0$ 时, $f(x)$ 是凹函数。

从定义不难证明: $f(x)(x\in I)$ 为凸函数的充分必要条件是对任意的 $x_1,x_2,x_3\in I$, 只要 $x_1<x_3<x_2$, 便有
\begin{equation}
  \frac{f(x_3)-f(x_1)}{x_3-x_1}
  \le
  \frac{f(x_2)-f(x_3)}{x_2-x_3}.
\tag{5.4.2}
\end{equation}
事实上, 取 $t=\dfrac{x_2-x_3}{x_2-x_1}$, 则 $x_3=tx_1+(1-t)x_2$, 再利用凸函数定义即得 (5.4.2) 式. 反之, 若 (5.4.2) 成立, 可将其化简为
\[
  f(x_3)\le
  \frac{x_2-x_3}{x_2-x_1}f(x_1)
  +\frac{x_3-x_1}{x_2-x_1}f(x_2).
\]
令 $t=\dfrac{x_2-x_3}{x_2-x_1}$, 则 $t\in(0,1),\ x_1<x_3=tx_1+(1-t)x_2\le x_2$, 则有
\[
  f(tx_1+(1-t)x_2)\le tf(x_1)+(1-t)f(x_2),
\]
所以 $f(x)$ 为凸函数。

当 $f(x)$ 为严格凸函数时, 不等式 (5.4.2) 成立严格不等式. 不等式 (5.4.2) 的几何解释是: 对一个凸函数而言, 过其图像上任意一点向两边作割线, 则左边割线的斜率总是小于或等于右边割线的斜率 (参见图 5.4.3)。

下面我们主要研究可导函数的凸凹性. 我们有:

\noindent\textbf{定理 5.4.4}\quad 设函数 $f(x)\in C[a,b]$ 且在 $(a,b)$ 上可导, 则 $f(x)$ 为凸函数的充分必要条件是 $f'(x)$ 在 $(a,b)$ 内单调上升; 而 $f(x)$ 为严格凸函数的充分必要条件是 $f'(x)$ 在 $(a,b)$ 内严格单调上升。

\noindent\textbf{证明}\quad 必要性\quad $\forall x_1,x_2\in(a,b),\ x_1<x_2$, 要证 $f'(x_1)\le f'(x_2)$. 令 $h>0$, 使得 $x_1-h,\ x_2+h\in(a,b)$, 则由凸函数的等价条件 (5.4.2), 有
\[
  \frac{f(x_1)-f(x_1-h)}{h}
  \le
  \frac{f(x_2)-f(x_1)}{x_2-x_1}
  \le
  \frac{f(x_2+h)-f(x_2)}{h}.
\]
令 $h\to0$, 得
\[
  f'(x_1)\le\frac{f(x_2)-f(x_1)}{x_2-x_1}\le f'(x_2),
\]
即 $f'(x)$ 在 $(a,b)$ 内单调上升。

若 $f(x)$ 为严格凸函数, 在 $(x_1,x_2)$ 中任取一点 $x^*$, 这时有
\[
  \frac{f(x_1)-f(x_1-h)}{h}
  <
  \frac{f(x^*)-f(x_1)}{x^*-x_1}
  <
  \frac{f(x_2)-f(x^*)}{x_2-x^*}
  <
  \frac{f(x_2+h)-f(x_2)}{h}.
\]
令 $h\to0$, 得
\[
  f'(x_1)\le
  \frac{f(x^*)-f(x_1)}{x^*-x_1}
  <
  \frac{f(x_2)-f(x^*)}{x_2-x^*}
  \le f'(x_2),
\]
即 $f'(x)$ 在 $(a,b)$ 内严格单调上升。

充分性\quad 要证 $f(x)$ 为凸函数, 由等价条件 (5.4.2) 知, 只需证对任意的 $x_1,x_2,x_3\in[a,b]$, 当 $x_1<x_2<x_3$ 时, 有
\[
  \frac{f(x_2)-f(x_1)}{x_2-x_1}
  \le
  \frac{f(x_3)-f(x_2)}{x_3-x_2}.
\]
由拉格朗日微分中值定理, 可得
\[
  \frac{f(x_2)-f(x_1)}{x_2-x_1}=f'(\xi_1),\qquad
  \frac{f(x_3)-f(x_2)}{x_3-x_2}=f'(\xi_2),
\]
其中 $x_1<\xi_1<x_2<\xi_2<x_3$. 由 $f'(x)$ 在 $(a,b)$ 内单调上升知, $f'(\xi_1)\le f'(\xi_2)$, 从而
\[
  \frac{f(x_2)-f(x_1)}{x_2-x_1}
  \le
  \frac{f(x_3)-f(x_2)}{x_3-x_2},
\]
即 $f(x)$ 为凸函数. 若 $f'(x)$ 严格单调上升, 则可导出上述不等式严格成立, 从而 $f(x)$ 为严格凸函数。

\noindent\textbf{定理 5.4.5}\quad 设函数 $f(x)\in C[a,b]$ 且在 $(a,b)$ 上二阶可导, 则 $f(x)$ 为凸函数的充分必要条件为 $f''(x)\ge0$; 而 $f(x)$ 为严格凸函数的充分必要条件为:

(1) $f''(x)\ge0$;

(2) $f''(x)$ 在 $(a,b)$ 的任一子区间上都不恒等于 $0$。

证明留作练习。

\noindent\textbf{例 5.4.8}\quad 设 $f(x)$ 是 $[a,b]$ 上的凸函数. 证明: 对任意的
\[
  x_1,x_2,\cdots,x_n\in[a,b],\qquad
  \sum_{i=1}^n t_i=1\quad (t_i>0),
\]
有
\[
  f(t_1x_1+\cdots+t_nx_n)\le t_1f(x_1)+\cdots+t_nf(x_n),
\]
而且当 $f(x)$ 为严格凸函数且 $x_i$ 不全相等时, 上述不等式严格成立。

\noindent\textbf{证明}\quad 用数学归纳法. 当 $n=2$ 时, 就是凸函数的定义. 现假定该命题对 $n=k$ 成立, 下证 $n=k+1$ 也成立。

对任意的 $x_i\in[a,b],\ t_i>0\ (i=1,2,\cdots,k+1),\ \sum_{i=1}^{k+1}t_i=1$, 令
\[
  \lambda_i=\frac{t_i}{1-t_{k+1}},\qquad i=1,2,\cdots,k,
\]
则有 $\lambda_i>0\ (i=1,2,\cdots,k)$, 而且 $\sum_{i=1}^k\lambda_i=1$. 于是, 由凸函数的定义和归纳法假设, 可得
\[
\begin{aligned}
  &f(t_1x_1+\cdots+t_kx_k+t_{k+1}x_{k+1})\\
  &\quad=f[(1-t_{k+1})(\lambda_1x_1+\cdots+\lambda_kx_k)+t_{k+1}x_{k+1}]\\
  &\quad\le(1-t_{k+1})f(\lambda_1x_1+\cdots+\lambda_kx_k)+t_{k+1}f(x_{k+1})\\
  &\quad\le(1-t_{k+1})(\lambda_1f(x_1)+\cdots+\lambda_kf(x_k))+t_{k+1}f(x_{k+1})\\
  &\quad=t_1f(x_1)+\cdots+t_kf(x_k)+t_{k+1}f(x_{k+1}).
\end{aligned}
\]
当 $f(x)$ 为严格凸函数且 $x_i$ 不全相等时, 分两种情况: 一种情况是 $x_1,x_2,\cdots,x_k$ 不全相等, 此时由归纳法假设, 可知严格不等式成立; 另一种情况是 $x_1,x_2,\cdots,x_k$ 相等, 但不等于 $x_{k+1}$, 此时由 $\lambda_1x_1+\cdots+\lambda_kx_k\ne x_{k+1}$, 可导出严格不等式也成立. 这样, 由归纳法原理知这一命题对一切正整数 $n$ 成立。

\noindent\textbf{例 5.4.9}\quad 设 $a_i>0\ (i=1,2,\cdots,n)$ 不全相等, 证明: 当 $x\ne0$ 时, 有
\[
  x\frac{a_1^x\ln a_1+\cdots+a_n^x\ln a_n}{a_1^x+\cdots+a_n^x}
  -\ln\frac{a_1^x+\cdots+a_n^x}{n}>0.
\]

\noindent\textbf{证明}\quad 简单的运算可知, 所要证的不等式等价于
\[
  \frac1n a_1^x\ln a_1^x+\cdots+\frac1n a_n^x\ln a_n^x
  >\frac{a_1^x+\cdots+a_n^x}{n}\ln\frac{a_1^x+\cdots+a_n^x}{n}.
\]
现在在区间 $(0,+\infty)$ 上考虑函数 $f(u)=u\ln u$, 则有
\[
  f'(u)=\ln u+1,\qquad f''(u)=\frac1u>0\quad (u>0).
\]
所以 $f(u)$ 是 $(0,+\infty)$ 上的严格凸函数. 又 $a_i^x\ (i=1,2,\cdots,n)$ 不全相等, 应用上例的结论, 得
\[
  f\left(\frac{a_1^x+\cdots+a_n^x}{n}\right)
  <\frac1n f(a_1^x)+\cdots+\frac1n f(a_n^x).
\]
这正是所要证明的不等式。

\noindent\textbf{例 5.4.10}\quad 设 $a_i>0\ (i=1,2,\cdots,n)$ 不全相等, 证明
\[
  f(x)=
  \begin{cases}
    \left(\dfrac{a_1^x+\cdots+a_n^x}{n}\right)^{\frac1x}, & x\ne0,\\[6pt]
    \sqrt[n]{a_1a_2\cdots a_n}, & x=0,
  \end{cases}
\]
是 $(-\infty,+\infty)$ 上的严格单调递增函数。

\noindent\textbf{证明}\quad 由例 5.2.8 可知, $f(x)\in C(-\infty,+\infty)$. 当 $x\ne0$ 时, 对 $\ln f(x)$ 求导, 得
\[
  \frac{f'(x)}{f(x)}
  =\frac1{x^2}\left(
  x\frac{a_1^x\ln a_1+\cdots+a_n^x\ln a_n}{a_1^x+\cdots+a_n^x}
  -\ln\frac{a_1^x+\cdots+a_n^x}{n}\right).
\]
因为 $f(x)>0$, 利用上例所证不等式可知, 当 $x\ne0$ 时, $f'(x)>0$. 因此, $f(x)$ 在实轴上是严格单调递增的。

\noindent\textbf{注}\quad 对上例的 $f(x)$, 我们有:
\[
  f(-1)=\frac{n}{\dfrac1{a_1}+\dfrac1{a_2}+\cdots+\dfrac1{a_n}},
\]
称之为 $a_1,a_2,\cdots,a_n$ 的调和平均;
\[
  f(0)=\sqrt[n]{a_1a_2\cdots a_n},
\]
称之为 $a_1,a_2,\cdots,a_n$ 的几何平均;
\[
  f(1)=\frac{a_1+a_2+\cdots+a_n}{n},
\]
称之为 $a_1,a_2,\cdots,a_n$ 的算术平均. 另外, 例 5.2.8 表明
\[
  f(-\infty)=\lim_{x\to-\infty}f(x)=\min\{a_1,a_2,\cdots,a_n\},
\]
\[
  f(+\infty)=\lim_{x\to+\infty}f(x)=\max\{a_1,a_2,\cdots,a_n\}.
\]
这样, 由 $f(x)$ 的严格单调递增性, 我们有
\[
  f(-\infty)<f(-1)<f(0)<f(1)<f(+\infty),
\]
即
\[
  \min\{a_1,a_2,\cdots,a_n\}
  <\frac{n}{\dfrac1{a_1}+\dfrac1{a_2}+\cdots+\dfrac1{a_n}}
  <\sqrt[n]{a_1a_2\cdots a_n}
\]
\[
  <\frac{a_1+a_2+\cdots+a_n}{n}
  <\max\{a_1,a_2,\cdots,a_n\}.
\]
这就是著名的调和--几何--算术平均不等式。

\noindent\textbf{例 5.4.11}\quad 设 $a_i>0,\ b_i>0\ (i=1,2,\cdots,n)$, 证明:
\[
  \sum_{i=1}^n a_i b_i
  \le
  \left(\sum_{i=1}^n a_i^p\right)^{\frac1p}
  \left(\sum_{i=1}^n b_i^q\right)^{\frac1q},
\]
其中 $0<p,q<+\infty,\ \dfrac1p+\dfrac1q=1$. 此不等式称为赫尔德 (H\"older) 不等式, 当 $p=q=2$ 时, 又称为施瓦茨 (Schwarz) 不等式或柯西不等式, 它表明 $n$ 维空间中的两个向量夹角余弦之绝对值不超过 1。

\noindent\textbf{证明}\quad 令 $f(x)=x^{\frac1q}$, 则
\[
  f''(x)=\frac1q\left(\frac1q-1\right)x^{\frac1q-2}<0,\qquad \forall x>0.
\]
因此, $f(x)$ 为 $(0,+\infty)$ 上的严格凹函数. 于是, 若
\[
  x_i>0,\qquad t_i>0,\qquad \sum_{i=1}^n t_i=1,
\]
则有
\[
  t_1x_1^{\frac1q}+t_2x_2^{\frac1q}+\cdots+t_nx_n^{\frac1q}
  \le (t_1x_1+t_2x_2+\cdots+t_nx_n)^{\frac1q}.
\]
现取
\[
  t_i=\frac{a_i^p}{\sum_{i=1}^n a_i^p},\qquad x_i=\frac{b_i^q}{a_i^p},
\]
并且代入上述不等式, 得
\[
  \frac{a_1b_1+\cdots+a_nb_n}{\sum_{i=1}^n a_i^p}
  \le
  \frac{(b_1^q+\cdots+b_n^q)^{\frac1q}}{\left(\sum_{i=1}^n a_i^p\right)^{\frac1q}},
\]
整理即得
\[
  \sum_{i=1}^n a_i b_i
  \le
  \left(\sum_{i=1}^n a_i^p\right)^{\frac1p}
  \left(\sum_{i=1}^n b_i^q\right)^{\frac1q}.
\]

\subsection{拐点}

\noindent\textbf{定义 5.4.2}\quad 设函数 $f(x)$ 在 $U(x_0,\delta)$ 上连续. 如果存在 $\delta_0>0$, 使得 $f(x)$ 在 $U_0^-(x_0,\delta_0)$ 是凹 (凸) 的, 而 $U_0^+(x_0,\delta_0)$ 是凸 (凹) 的, 则称 $x_0$ 为 $f(x)$ 的一个拐点。

\noindent\textbf{定理 5.4.6}\quad 如果 $x_0$ 是 $f(x)$ 的拐点, 且 $f''(x_0)$ 存在, 则
\[
  f''(x_0)=0.
\]

\noindent\textbf{证明}\quad $f''(x_0)$ 存在表明 $f'(x)$ 在点 $x_0$ 附近存在. $f(x)$ 在点 $x_0$ 的左右凹凸性相反, 表明 $f'(x)$ 在点 $x_0$ 的左右升降性相反, 即 $x_0$ 是 $f'(x)$ 一个极值点. 由费马定理知, $f''(x_0)=0$。

\noindent\textbf{定理 5.4.7}\quad 如果函数 $f(x)$ 在 $U(x_0,\delta)$ 二阶可导, $f''(x_0)=0$, $f'''(x_0)$ 存在而且不为零, 则 $x_0$ 是 $f(x)$ 的拐点。

\noindent\textbf{证明}\quad $f''(x)$ 的带佩亚诺余项的泰勒公式为
\[
  f''(x)=f''(x_0)+f'''(x_0)(x-x_0)+o(x-x_0)
  =(f'''(x_0)+o(1))(x-x_0).
\]
所以 $\exists\delta_1<\delta$, 使得当 $x\in U(x_0,\delta_1)$ 时, $f'''(x_0)+o(1)$ 与 $f'''(x_0)$ 有相同符号, 从而 $f''(x)$ 在 $x_0$ 左右符号相反, 即 $f(x)$ 在 $x_0$ 左右凸凹性相反, 故 $x_0$ 是拐点。

从拐点定义容易看出, 若函数 $f(x)$ 具有二阶导数, 则 $f(x)$ 的拐点即为 $f'(x)$ 的极值点. 因此将前面关于 $f(x)$ 的极值点的结果应用于 $f'(x)$, 即得到了关于拐点的相应结果。

\subsection{渐近线}

设函数 $y=f(x)$ 的图像有一向无穷远无限伸展的分支, 如果当一个点沿着这一分支趋向无穷远时, 该点与某一条直线 $\ell$ 的距离趋向于零, 则称直线 $\ell$ 为 $f(x)$ 的一条渐近线。

首先, 如果有 $\lim_{x\to x_0}f(x)=\infty$, 则 $x=x_0$ 就是 $f(x)$ 的一条渐近线. 此时, 我们称它为 $f(x)$ 的一条垂直渐近线. 当然也可以考虑 $x\to x_0+0$ 和 $x\to x_0-0$ 的情形。

下面考虑非垂直的渐近线. 设 $Y=kx+b$ 为 $y=f(x)$ 的一条渐近线 (参见图 5.4.4). 由解析几何的知识知, 函数 $y=f(x)$ 的图像上一点 $M=(x,f(x))$ 到直线 $Y=kx+b$ 的距离为
\[
  \delta=|y-Y|\cos\alpha=\frac{|f(x)-kx-b|}{\sqrt{1+k^2}}.
\]
因此, 当这点沿着曲线 $y=f(x)$ 趋于无穷时, 它与直线 $Y=kx+b$ 的距离 $\delta$ 趋于零的充要条件是
\[
  \lim_{x\to+\infty}[f(x)-(kx+b)]=0.
\]
于是
\[
  \lim_{x\to+\infty}\frac{f(x)}x
  =\lim_{x\to+\infty}\left(\frac{f(x)-kx-b}{x}+\frac{kx+b}{x}\right)=k.
\]
这样, 所求渐近线的斜率 $k$ 和截距 $b$ 为
\[
  k=\lim_{x\to+\infty}\frac{f(x)}x,\qquad
  b=\lim_{x\to+\infty}(f(x)-kx).
\]
应当注意, 非垂直渐近线包括了斜渐近线和水平渐近线. 当由上面的式子确定的 $k=0$ 时, 则有
\[
  b=\lim_{x\to+\infty}f(x).
\]
此时, $y=b$ 就是一条水平渐近线。

同理可以考虑 $x\to-\infty$ 和 $x\to\infty$ 的情形。

\subsection{函数的作图}

解决实际问题时, 遇到的总是具体函数, 作出函数的图像, 对我们分析问题和解决问题会有很大的帮助。

最基本的作图方法是描点画图法, 计算机就是利用这个方法来作图的. 一般来说, 我们先把 $[a,b]$ 分得充分细, 在每个分点 $x_i$ 上计算 $f(x_i)$; 然后描出点 $(x_i,f(x_i))$, 并且用光滑的曲线将各点连起来, 便得到了 $y=f(x)$ 的图像. 本书中的图全部是这样做出来的。

手工作图不能这样, 因为计算量太大. 但利用导数作为工具, 先把函数各种性质尽可能地搞清楚, 就可很快作出函数的大致图像. 其具体步骤如下:

(1) 求出函数的定义域;

(2) 研究函数的有界性、奇偶性和周期性;

(3) 解方程 $f'(x)=0$, 列表确定函数升降区间和极值点;

(4) 解方程 $f''(x)=0$, 列表确定函数的凸凹区间和拐点;

(5) 求出函数的斜渐近线与垂直渐近线;

(6) 计算一些重要点上 (如点 $x=0$) 的函数值;

(7) 根据以上数据及函数的变化趋势描出其草图。

\noindent\textbf{例 5.4.12}\quad 作出函数 $y=e^{-x^2}$ 的草图。

\noindent\textbf{解}\quad 显然, 函数 $y=e^{-x^2}$ 定义域为 $(-\infty,+\infty)$, 且为偶函数. 因此它的图像关于 $y$ 轴对称. 当 $x=0$ 时, $y=1$, 并且总有 $y>0$, 可见它的图像仅与 $y$ 轴相交而与 $x$ 轴不相交. 此外, 由于
\[
  \lim_{x\to\pm\infty}e^{-x^2}=0,
\]
所以 $x$ 轴为它的一条水平渐近线. 求此函数的导数并且解方程 $y'=-2xe^{-x^2}=0$, 知其有唯一的驻点 $x=0$. 再求此函数的二阶导数并且解方程
\[
  y''=4e^{-x^2}\left(x^2-\frac12\right)=0,
\]
得 $x=\pm1/\sqrt2$. 关于这一函数的升降性、极值、凹凸性及拐点的情况, 如表 5.4.1 所示。

根据上面的分析, 大致描出 $y=e^{-x^2}$ 的图像如图 5.4.5 所示。

\noindent\textbf{例 5.4.13}\quad 作出函数
\[
  f(x)=\frac{(x-1)^3}{(x+1)^2}
\]
的草图。

\noindent\textbf{解}\quad 除 $x=-1$ 外, $f(x)$ 在实轴上都有意义. 显然有 $\lim_{x\to-1}f(x)=-\infty$, 因此 $x=-1$ 是该函数的一条垂直渐近线. 又
\[
  \lim_{x\to\infty}\frac{f(x)}x=1,\qquad
  \lim_{x\to\infty}(f(x)-x)=-5,
\]
所以 $y=x-5$ 是它的一条斜渐近线. 此外, 直接计算有 $f(1)=0$ 及
\[
  f'(x)=\frac{(x-1)^2(x+5)}{(x+1)^3},\qquad
  f'(1)=0,\qquad f'(-5)=0,
\]
\[
  f''(x)=\frac{24(x-1)}{(x+1)^4},\qquad f''(1)=0.
\]
根据上面的结果, 可得这一函数的升降性、极值、凹凸性及拐点的情况, 如表 5.4.2 所示。

根据上面的分析, 大致描出该函数的图像如图 5.4.6 所示。

\section*{习题五}

1. 证明广义罗尔微分中值定理: 设函数 $f(x)$ 在 $(a,b)$ 上可导, $f(a+0)=f(b-0)=A$, 则存在 $\xi\in(a,b)$, 使得 $f'(\xi)=0$, 其中 $a$ 可为 $-\infty$, $b$ 可为 $+\infty$, $A$ 可为 $+\infty$ 或 $-\infty$。

2. 证明: 函数 $f(x)=[(x-a)^n(x-b)^n]^{(n)}$ 在 $(a,b)$ 内有 $n$ 个互不相同的零点。

3. 证明: 若函数 $f(x)=x^m(1-x)^n$, 其中 $m,n$ 为正整数, 则存在 $\xi\in(0,1)$, 使得
\[
  \frac mn=\frac{\xi}{1-\xi}.
\]

4. 求证: $4ax^3+3bx^2+2cx=a+b+c$ 在 $(0,1)$ 上至少有一个根。

5. 求证: 切比雪夫--拉盖尔 (Chebyshev-Laguerre) 多项式
\[
  L_n(x)=e^x\frac{\mathrm d^n}{\mathrm dx^n}(x^ne^{-x})
\]
有 $n$ 个不同的零点。

6. 求证: 切比雪夫--埃尔米特 (Chebyshev-Hermite) 多项式
\[
  H_n(x)=(-1)^n\frac1{n!}e^{\frac{x^2}2}\frac{\mathrm d^n}{\mathrm dx^n}\left(e^{-\frac{x^2}2}\right)
\]
有 $n$ 个不同的零点。

7. 证明:

(1) 若函数 $f(x)$ 在 $(a,b)$ 上的导数有界, 则 $f(x)$ 在 $(a,b)$ 上一致连续;

(2) 对任意的两个实数 $x_1,x_2$, 有
\[
  |\sin x_1-\sin x_2|\le |x_1-x_2|;
\]

(3) 对任意的两个实数 $x_1,x_2$, 有
\[
  |\arctan x_1-\arctan x_2|\le |x_1-x_2|.
\]

8. 设 $\lim_{x\to+\infty}f'(x)=a$, 求证: 对任意的 $T>0$, 有
\[
  \lim_{x\to+\infty}(f(x+T)-f(x))=Ta.
\]

9. 证明下列不等式:

(1) $x-\dfrac{x^3}6<\sin x<x\quad (x>0)$;

(2) $x-\dfrac{x^2}2<\ln(1+x)<x\quad (x>0)$;

(3) $2\sqrt x>3-\dfrac1x\quad (x>1)$;

(4) $e^x-1>x+\dfrac{x^2}2\quad (x>0)$;

(5) $\ln(1+x)>\dfrac{\arctan x}{1+x}\quad (x>0)$。

10. 设函数 $f(x)$ 和 $g(x)$ 在 $(-\infty,+\infty)$ 上可导, 且有
\[
  f'(x)>g'(x),\qquad f(a)=g(a).
\]
证明: 当 $x>a$ 时, 有 $f(x)>g(x)$; 而当 $x<a$ 时, 有 $f(x)<g(x)$。

11. 设函数 $f(x)$ 在区间 $[a,b]\ (0<a<b)$ 上连续, 在 $(a,b)$ 内可导. 求证: 存在 $\xi_1,\xi_2,\xi_3\in(a,b)$, 使得
\[
  f'(\xi_1)=(b+a)\frac{f'(\xi_2)}{2\xi_2}
  =(b^2+ab+a^2)\frac{f'(\xi_3)}{3\xi_3^2}.
\]

12. 设函数 $f(x)$ 在 $(a,b)$ 上可导, 且 $f'(x)$ 单调. 证明 $f'(x)$ 在 $(a,b)$ 上连续。

13. 设函数 $f(x)$ 在点 $a$ 二阶可导, 求证:
\[
  \lim_{h\to0}\frac{f(a+h)+f(a-h)-2f(a)}{h^2}=f''(a).
\]

14. 设函数 $f(x)$ 在 $[a,b]$ 上二阶可导, $f(a)=f(b)=0$, 且存在 $c\in(a,b)$, 使得 $f(c)>0$. 求证: 存在 $\xi\in(a,b)$, 使得 $f''(\xi)<0$。

15. 证明: 设 $x\ge0$, 则

(1)
\[
  \sqrt{x+1}-\sqrt x=\frac1{2\sqrt{x+\theta(x)}},\qquad
  \frac14\le\theta(x)\le\frac12;
\]

(2)
\[
  \lim_{x\to+0}\theta(x)=\frac14,\qquad
  \lim_{x\to+\infty}\theta(x)=\frac12.
\]

16. 设函数 $f(x)$ 在 $[a,b]$ 上三阶可导, 证明: 存在一点 $\xi\in(a,b)$, 使得
\[
  f(b)=f(a)+\frac12(b-a)[f'(a)+f'(b)]-\frac1{12}(b-a)^3f^{(3)}(\xi).
\]

17. 设函数
\[
  f(x)=
  \begin{cases}
    x^2\sin\dfrac1x, & x\ne0,\\[4pt]
    0, & x=0,
  \end{cases}
\]
在 $[0,x]$ 上应用中值定理, 得
\[
  x^2\sin\frac1x=x\left(2\xi\sin\frac1\xi-\cos\frac1\xi\right),\qquad 0<\xi<x.
\]
由此得
\[
  \cos\frac1\xi=2\xi\sin\frac1\xi-x\sin\frac1x.
\]
当 $x\to0$ 时, 有 $\xi\to0$, 所以由上式得到
\[
  \lim_{\xi\to0}\cos\frac1\xi=0.
\]
但已知 $\lim_{\xi\to0}\cos\dfrac1\xi$ 不存在. 试说明矛盾何在。

18. 设函数 $f(x)$ 在 $(0,+\infty)$ 上可导, 且 $\lim_{x\to+\infty}f'(x)=+\infty$. 求证 $f(x)$ 在 $(0,+\infty)$ 上不一致连续。

19. 证明函数 $f(x)=x\ln x$ 在 $(0,+\infty)$ 上不一致连续。

20. 设函数 $f(x),g(x),h(x)$ 都在 $[a,b]$ 上连续, 在 $(a,b)$ 内可导, 并定义
\[
  F(x)=
  \begin{vmatrix}
    f(a) & g(a) & h(a)\\
    f(b) & g(b) & h(b)\\
    f(x) & g(x) & h(x)
  \end{vmatrix}.
\]
证明: 存在 $\xi\in(a,b)$, 使得 $F'(\xi)=0$. 问若 $h(x)\equiv1$, 则上述结论与柯西中值公式有什么关系? 若令 $g(x)=x,\ h(x)\equiv1$, 则会有什么结论?

21. 设函数 $f(x)$ 在 $[a,b]$ 上连续, 在 $(a,b)$ 内可导. 求证: 存在 $\xi\in(a,b)$, 使得
\[
  2\xi(f(b)-f(a))=(b^2-a^2)f'(\xi).
\]

22. 设函数 $f(x)$ 在 $[a,b]$ 上连续, 在 $(a,b)$ 内可导. 求证: 存在 $\xi\in(a,b)$, 使得
\[
  f(b)-f(a)=\xi\ln\frac ba f'(\xi).
\]

23. 设函数 $f(x)$ 在 $[a,b]$ 上连续, 在 $(a,b)$ 内可导, 且 $f(a)=f(b)=1$. 求证: 存在 $\xi,\eta\in(a,b)$, 使得
\[
  e^{\eta-\xi}[f(\eta)+f'(\eta)]=1.
\]

24. 证明下列恒等式:

(1)
\[
  2\arctan x+\arcsin\frac{2x}{1+x^2}=\pi\operatorname{sgn}x\quad (|x|\ge1);
\]

(2)
\[
  3\arccos x-\arccos(3x-4x^3)=\pi\quad (|x|\le1/2).
\]

25. 求下列极限:

(1) $\displaystyle\lim_{x\to1}\frac{\ln\cos(x-1)}{1-\sin\frac{\pi x}2}$;\qquad
(2) $\displaystyle\lim_{x\to+\infty}(\pi-2\arctan x)\ln x$;

(3) $\displaystyle\lim_{x\to1}\left(\frac1{\ln x}-\frac1{x-1}\right)$;\qquad
(4) $\displaystyle\lim_{x\to1}\frac{x-1}{\ln x}$;

(5) $\displaystyle\lim_{x\to+0}x^{\sin x}$;\qquad
(6) $\displaystyle\lim_{x\to0}\left(\frac{(\ln(1+x))^{1+x}}{x^2}-\frac1x\right)$;

(7) $\displaystyle\lim_{x\to0}\left(\cot x-\frac1x\right)$;\qquad
(8) $\displaystyle\lim_{x\to0}\frac{(1+x)^{\frac1x}-e}{x}$;

(9) $\displaystyle\lim_{x\to+\infty}\left(\frac\pi2-\arctan x\right)^{\frac1{\ln x}}$;\qquad
(10) $\displaystyle\lim_{x\to+0}\sin x\ln x$;

(11) $\displaystyle\lim_{x\to0}(1-\cos x)^{1-\cos x}$;\qquad
(12) $\displaystyle\lim_{x\to0}\frac{\tan x-\sin x}{\sin x-x\cos x}$;

(13) $\displaystyle\lim_{x\to1}x^{\frac1{1-x}}$;\qquad
(14) $\displaystyle\lim_{x\to0}\frac{xe^x-\ln(1+x)}{x^2}$;

(15) $\displaystyle\lim_{x\to0}\frac{x^2\sin\dfrac1x}{\sin x}$;\qquad
(16) $\displaystyle\lim_{x\to+0}\left(\frac{1+x^a}{1+x^b}\right)^{\frac1{\ln x}}\quad (a,b\text{ 为正实数})$;

(17) $\displaystyle\lim_{x\to\frac\pi4-0}(\tan x)^{\tan 2x}$。

26. 设函数 $f(x)$ 在 $[a,+\infty)$ 上有界, $f'(x)$ 存在且 $\lim_{x\to+\infty}f'(x)=b$. 求证: $b=0$。

27. 由拉格朗日微分中值定理, 有
\[
  \ln(1+x)=\frac{x}{1+\theta x}\qquad (0<\theta<1).
\]
求证 $\lim_{x\to0}\theta=\dfrac12$。

28. 由拉格朗日微分中值定理, 有
\[
  \arcsin x=\frac{x}{\sqrt{1-\theta^2x^2}}\qquad (0<\theta<1).
\]
求证 $\lim_{x\to0}\theta=\dfrac1{\sqrt3}$。

29. 设函数 $f(x)$ 在 $(a,+\infty)$ 上可导, 且 $\lim_{x\to+\infty}[f(x)+f'(x)]=k$ ($k$ 有限或 $\pm\infty$). 求证: $\lim_{x\to+\infty}f(x)=k$。

30. 设函数 $f(x)$ 在实轴上任意阶可导, 令 $F(x)=f(x^2)$, 求证:
\[
  F^{(2n+1)}(0)=0,\qquad
  \frac{F^{(2n)}(0)}{(2n)!}=\frac{f^{(n)}(0)}{n!}.
\]

31. 写出下列函数在 $x=0$ 的带佩亚诺余项的泰勒公式:

(1) $\cos x^2$;\qquad
(2) $\sin^3x$;

(3) $\dfrac1{(1+x)^2}$;\qquad
(4) $\dfrac{x^3+2x+1}{x-1}$。

32. 写出下列函数在 $x=0$ 的带有佩亚诺余项的泰勒公式, 要求至所指定的阶数。

(1) $\dfrac{x}{\sin x}\quad (x^4)$;\qquad
(2) $\ln(\cos x+\sin x)\quad (x^4)$;

(3) $\dfrac{x}{2x^2+x-1}\quad (x^3)$;\qquad
(4) $\ln\dfrac{1+x}{1-2x}\quad (x^n)$;

(5) $\ln(1+x+x^2+x^3)\quad (x^6)$;\qquad
(6) $\dfrac{1+x+x^2}{1-x+x^2}\quad (x^4)$。

33. 确定常数 $a,b$, 使当 $x\to0$ 时,

(1) $f(x)=(a+b\cos x)\sin x-x$ 为 $x$ 的 5 阶无穷小量;

(2) $f(x)=e^x-\dfrac{1+ax}{1-bx}$ 是 $x$ 的 3 阶无穷小量。

34. 求下列极限:

(1) $\displaystyle\lim_{x\to0}\left(\frac1x-\frac1{\sin x}\right)$;\qquad
(2) $\displaystyle\lim_{x\to0}\frac{e^{x^3}-1-x^3}{\sin^6 2x}$;

(3) $\displaystyle\lim_{x\to+\infty}\left(\sqrt[3]{x^3-3x}-\sqrt{x^2-2x}\right)$;\qquad
(4) $\displaystyle\lim_{x\to\infty}\left(x+\frac12\right)\ln\left(1+\frac1x\right)$;

(5) $\displaystyle\lim_{x\to0}\left(\frac{\tan x}{x}\right)^{\frac1{x^2}}$;\qquad
(6) $\displaystyle\lim_{x\to\infty}x^2\ln\left(x\sin\frac1x\right)$。

35. 设函数 $f(x)$ 在 $x=0$ 的某个邻域内二阶可导, 且
\[
  \lim_{x\to0}\left(\frac{\sin3x}{x^3}+\frac{f(x)}{x^2}\right)=0.
\]

(1) 求 $f(0),f'(0),f''(0)$;

(2) 求 $\displaystyle\lim_{x\to0}\left(\frac3{x^2}+\frac{f(x)}{x^2}\right)$。

36. 设函数 $f(x)$ 在 $x=0$ 的某个邻域内二阶可导, 且
\[
  \lim_{x\to0}\left(1+x+\frac{f(x)}x\right)^{\frac1x}=e^3.
\]

(1) 求 $f(0),f'(0),f''(0)$;

(2) 求 $\displaystyle\lim_{x\to0}\left(1+\frac{f(x)}x\right)^{\frac1x}$。

37. 设函数 $f(x)$ 在 $(a,+\infty)$ 上有直到 $n$ 阶导数, 且有
\[
  \lim_{x\to+\infty}f(x)=A,\qquad
  \lim_{x\to+\infty}f^{(n)}(x)=B.
\]
求证 $B=0$。

38. 设函数 $f(x)$ 在 $x=a$ 的某个邻域内二阶导数连续, 且 $f''(a)\ne0$, 由拉格朗日微分中值定理有
\[
  f(a+h)-f(a)=f'(a+\theta h)h\qquad (0<\theta<1).
\]
求证:
\[
  \lim_{h\to0}\theta=\frac12.
\]

39. 设 $P(x)$ 为 $n$ 次多项式, 证明:

(1) 若 $P(a),P'(a),\cdots,P^{(n)}(a)$ 皆为正数, 则 $P(x)$ 在 $(a,+\infty)$ 上无零点;

(2) 若 $P(a),P'(a),\cdots,P^{(n)}(a)$ 正负号相间, 则 $P(x)$ 在 $(-\infty,a)$ 上无零点;

40. 设 $P(x)$ 为 $n$ 次多项式, 证明 $x=a$ 是 $P(x)=0$ 的 $k$ 重根的充分必要条件是
\[
  P^{(j)}(a)=0\ (j=0,1,\cdots,k-1),\qquad P^{(k)}(a)\ne0.
\]

41. 设
\[
  P_n(x)=\frac1{2^n n!}\frac{\mathrm d^n}{\mathrm dx^n}(x^2-1)^n,
\]
求证:

(1) $P_n(1)=1$;

(2) $P_n(-1)=(-1)^n$;

(3) $P_{2n-1}(0)=0$;

(4)
\[
  P_{2n}(0)=\frac{(-1)^n(2n)!}{2^{2n}(n!)^2}.
\]

42. 设函数 $f(x)=e^{x^2}$。

(1) 求证: $f^{(n)}(x)=P_n(x)e^{x^2}$, 其中 $P_n(x)$ 为 $n$ 次多项式, 且满足 $P_0(x)=1,\ P_1(x)=2x,\ P_{n+1}(x)=2xP_n(x)+2nP_{n-1}(x)$;

(2) 求 $f^{(n)}(0)$ 的值。

43. 设函数 $f(x)$ 在 $(0,+\infty)$ 上三阶可导, 而且有
\[
  |f(x)|\le M_0,\qquad |f'''(x)|\le M_3,\qquad \forall x\in(0,+\infty).
\]
求证 $f'(x)$ 和 $f''(x)$ 在 $(0,+\infty)$ 上有界。

44. 求证:

(1) $0<x-\ln(1+x)<\dfrac12x^2\quad (0<x\le1)$;

(2) $\displaystyle\lim_{n\to\infty}\sum_{k=1}^n\left[\frac1k-\ln\left(1+\frac1k\right)\right]$ 存在。

45. 设函数 $f(x)$ 在 $[a,b]$ 上有二阶导数, 且 $f'(a)=f'(b)=0$. 证明: 存在 $c\in(a,b)$, 使得
\[
  |f''(c)|\ge\frac4{(b-a)^2}|f(b)-f(a)|.
\]

46. 若函数 $f(x)$ 在 $[-1,1]$ 上三阶可导, 且有
\[
  f(0)=f'(0)=0,\qquad f(1)=1,\qquad f(-1)=0.
\]
求证: 存在 $\xi\in(-1,1)$, 使得 $f'''(\xi)\ge3$。

47. 设函数 $f(x)$ 在 $(-\infty,+\infty)$ 上二阶可导, 且有
\[
  |f(x)|\le M_0,\qquad |f''(x)|\le M_2,\qquad \forall x\in(-\infty,+\infty).
\]

(1) 写出 $f(x+h)$ 和 $f(x-h)$ 的泰勒公式;

(2) 求证: $|f'(x)|\le\dfrac{M_0}h+\dfrac h2M_2,\ \forall h>0$;

(3) 求 $\dfrac{M_0}h+\dfrac h2M_2$ 在 $(0,+\infty)$ 上的最小值;

\end{document}
